% Luc Destombe --- licence CC BY-NC-SA 4.0
% https://creativecommons.org/licenses/by-nc-sa/4.0/deed.fr
% Fichier autonome : compiler avec pdflatex, deux fois (signets, nombre de pages).
\def\lyceeversion{1.1}
\documentclass[11pt,a4paper]{article}
\makeatletter
\newif\ifeleve
\elevefalse

\newif\iflycee@cm
\newif\iflycee@palatino

\RequirePackage[utf8]{inputenc}
\RequirePackage[T1]{fontenc}
\RequirePackage[french]{babel}
\RequirePackage{etoolbox}

\newcommand{\ShorthandsOff}{\@ifundefined{shorthandoff}{}{\shorthandoff{;:!?}}}
\newcommand{\ShorthandsOn}{\@ifundefined{shorthandon}{}{\shorthandon{;:!?}}}
\ShorthandsOff
\AtBeginDocument{\ShorthandsOn}
\AtBeginEnvironment{tikzpicture}{\ShorthandsOff}

\RequirePackage{geometry}
\RequirePackage{fancyhdr}
\RequirePackage{lastpage}
\RequirePackage{multicol}
\RequirePackage{array}
\RequirePackage{enumitem}
\RequirePackage{xcolor}
\RequirePackage{needspace}

\RequirePackage{amsmath,amssymb}
\iflycee@cm\else\RequirePackage{fourier}\fi
\DeclareMathAlphabet{\mathcal}{OMS}{cmsy}{m}{n}
\renewcommand{\ttdefault}{lmtt}
\begingroup\catcode`\_=\active \catcode`\^=\active
\global\let_\sb \global\let^\sp
\endgroup
\AtBeginDocument{\catcode`\_=12 \mathcode`\_="8000
  \catcode`\^=12 \mathcode`\^="8000 }
\newlength{\retraitcalculs}
\setlength{\retraitcalculs}{2em}

\RequirePackage{tikz}
\usetikzlibrary{arrows.meta,calc,babel,shapes.misc}
\RequirePackage[most]{tcolorbox}
\tcbuselibrary{skins,breakable}

\definecolor{coulDef}{HTML}{1F4E79}
\definecolor{coulProp}{HTML}{7030A0}
\definecolor{coulRem}{HTML}{C55A11}
\definecolor{coulEx}{HTML}{2E7D32}
\definecolor{coulExo}{HTML}{B03A2E}
\definecolor{coulPart}{HTML}{1F4E79}
\definecolor{coulMeth}{HTML}{1B6E4F}
\definecolor{coulCourbe}{HTML}{0B5ED7}
\definecolor{coulCourbeB}{HTML}{C00000}
\definecolor{coulCourbeC}{HTML}{2E7D32}
\definecolor{coulGrille}{HTML}{8C8C8C}
\definecolor{coulRep}{HTML}{1B6E4F}
\definecolor{coulBar}{HTML}{2E7D32}

\newcommand{\R}{\mathbb{R}}
\newcommand{\Rs}{\mathbb{R}^{*}}
\newcommand{\N}{\mathbb{N}}
\newcommand{\Ns}{\mathbb{N}^{*}}
\newcommand{\Z}{\mathbb{Z}}
\newcommand{\Q}{\mathbb{Q}}
\renewcommand{\le}{\leqslant}
\renewcommand{\ge}{\geqslant}
\renewcommand{\approx}{\simeq}
\newcommand{\vect}[1]{\overrightarrow{#1}}
\newcommand{\norme}[1]{\left\lVert #1 \right\rVert}
\newcommand{\intervalle}[2]{\left[#1\,;#2\right]}
\RequirePackage{cancel}
\renewcommand{\CancelColor}{\color{coulExo}}
\newcommand{\fois}[1]{\times{\color{coulExo}#1}}
\newcommand{\barre}[1]{\cancel{#1}}
\newlength{\lycee@eqh}\newlength{\lycee@eqd}
\newsavebox{\lycee@eqa}\newsavebox{\lycee@eqb}\newsavebox{\lycee@eqc}
\newcommand{\lycee@eqcol}[1]{\hbox{$\renewcommand{\arraystretch}{1.35}%
  \begin{array}[t]{@{}>{\displaystyle}r@{}>{\displaystyle{}}l@{}}#1\end{array}$}}
\newcommand{\lycee@eqmesure}[1]{%
  \ifdim\ht#1>\lycee@eqh\setlength{\lycee@eqh}{\ht#1}\fi
  \ifdim\dp#1>\lycee@eqd\setlength{\lycee@eqd}{\dp#1}\fi}
\newcommand{\lycee@eqtrait}{\hspace{1.2em}\tikz[baseline=0pt]\draw[dash pattern=on 2pt off 2pt,
  line width=0.5pt] (0,\lycee@eqh) -- (0,-\lycee@eqd);\hspace{1.2em}}
\newcommand{\eqproduit}[2]{\par\smallskip\noindent C'est une équation produit :
  $A\times B=0 \iff A=0$ ou $B=0$.\par\nopagebreak\smallskip
  \begingroup\setlength{\lycee@eqh}{0pt}\setlength{\lycee@eqd}{0pt}%
  \setbox\lycee@eqa=\lycee@eqcol{#1}\setbox\lycee@eqb=\lycee@eqcol{#2}%
  \lycee@eqmesure\lycee@eqa\lycee@eqmesure\lycee@eqb%
  \noindent\hspace*{\retraitcalculs}\usebox\lycee@eqa\lycee@eqtrait\usebox\lycee@eqb\endgroup\par\smallskip}
\newcommand{\eqproduittrois}[3]{\par\smallskip\noindent C'est une équation produit :
  $A\times B\times C=0 \iff A=0$ ou $B=0$ ou $C=0$.\par\nopagebreak\smallskip
  \begingroup\setlength{\lycee@eqh}{0pt}\setlength{\lycee@eqd}{0pt}%
  \setbox\lycee@eqa=\lycee@eqcol{#1}\setbox\lycee@eqb=\lycee@eqcol{#2}\setbox\lycee@eqc=\lycee@eqcol{#3}%
  \lycee@eqmesure\lycee@eqa\lycee@eqmesure\lycee@eqb\lycee@eqmesure\lycee@eqc%
  \noindent\hspace*{\retraitcalculs}\usebox\lycee@eqa\lycee@eqtrait\usebox\lycee@eqb\lycee@eqtrait\usebox\lycee@eqc\endgroup
  \par\smallskip}

\tikzset{grille/.style={coulGrille,line width=0.4pt}}
\newcommand{\quadrillage}[4]{\draw[grille,step=1] (#1,#2) grid (#3,#4);}
\tikzset{
  croix/.style={cross out,draw,line width=0.6pt,inner sep=0pt,minimum size=4pt},
  croixdroite/.style={croix,rotate=45,line width=0.8pt},
}
\newcommand{\pt}[3]{\path (#1) node[croix]{} node[#2,font=\small,inner sep=2.5pt]{$#3$};}

\newcommand{\lycee@repere}[4]{%
  \draw[grille,step=1] (#1,#2) grid (#3,#4);
  \draw[-{Stealth[length=2mm]},thick] (#1,0)--({#3+0.4},0) node[below left=-1pt]{$x$};
  \draw[-{Stealth[length=2mm]},thick] (0,#2)--(0,{#4+0.4}) node[below left=-1pt]{$y$};
  \node[below left,font=\scriptsize,inner sep=1.5pt] at (0,0) {$0$};}
\newcommand{\lycee@gradx}[1]{\foreach \k/\l in {#1}{%
  \draw (\k,0.07)--(\k,-0.07) node[below,font=\scriptsize,inner sep=1.5pt]{$\l$};}}
\newcommand{\lycee@grady}[1]{\foreach \k/\l in {#1}{%
  \draw (0.07,\k)--(-0.07,\k) node[left,font=\scriptsize,inner sep=1.5pt]{$\l$};}}
\newcommand{\lycee@intff}[2]{\left[#1\,;#2\right]}
\newcommand{\lycee@intfo}[2]{\left[#1\,;#2\right[}
\newcommand{\lycee@intof}[2]{\left]#1\,;#2\right]}
\newcommand{\lycee@intoo}[2]{\left]#1\,;#2\right[}
\newcommand{\lycee@ens}[1]{\left\{#1\right\}}
\AtBeginDocument{%
  \providecommand{\repere}{\lycee@repere}%
  \providecommand{\gradx}{\lycee@gradx}%
  \providecommand{\grady}{\lycee@grady}%
  \providecommand{\Cf}{\mathcal{C}\sb{f}}%
  \providecommand{\Cg}{\mathcal{C}\sb{g}}%
  \providecommand{\intff}{\lycee@intff}%
  \providecommand{\intfo}{\lycee@intfo}%
  \providecommand{\intof}{\lycee@intof}%
  \providecommand{\intoo}{\lycee@intoo}%
  \providecommand{\ens}{\lycee@ens}}

\newlist{questions}{enumerate}{3}
\setlist[questions,1]{label=\arabic*\textdegree),leftmargin=*,itemsep=2pt,topsep=3pt}
\setlist[questions,2]{label=\alph*),leftmargin=*,itemsep=1pt,topsep=2pt}
\setlist[questions,3]{label=\roman*),leftmargin=*,itemsep=1pt,topsep=1pt}
\setlist[itemize]{label=\textbullet,leftmargin=*,itemsep=2pt,topsep=3pt}

\newcommand{\besoinplace}[1]{\par\penalty\@M\begingroup
  \dimen@=#1\relax\dimen@ii=\pagegoal\advance\dimen@ii-\pagetotal
  \ifdim\dimen@>\dimen@ii\ifdim\dimen@ii>\z@\vfil\fi\break\fi\endgroup}

\newcommand{\lycee@pied}{\thepage/\pageref{LastPage}}
\newcommand{\entete}[2]{%
  \pagestyle{fancy}\fancyhf{}%
  \fancyhead[L]{\footnotesize\color{coulPart}\textbf{#1}}%
  \fancyhead[R]{\footnotesize\color{coulPart}\textbf{#2}}%
  \fancyfoot[C]{\footnotesize\lycee@pied}%
  \renewcommand{\headrulewidth}{0.4pt}%
  \renewcommand{\headrule}{\hbox to\headwidth{%
    \color{coulPart!40}\leaders\hrule height \headrulewidth\hfill}}}

\RequirePackage{ccicons}
\newcommand{\lycee@licence}{{\scriptsize\color{gray}%
  \copyright\ Luc Destombe --- \ccbyncsa\ CC BY-NC-SA 4.0}}
\newif\iflycee@licence \lycee@licencetrue
\newcommand{\sanslicence}{\lycee@licencefalse}
\AtBeginDocument{\iflycee@licence\AddToHookNext{shipout/foreground}{%
  \put(0,0){\rlap{\kern\dimexpr1in+\hoffset+\oddsidemargin+\textwidth\relax
    \llap{\raisebox{-\dimexpr1in+\voffset+\topmargin+\headheight+\headsep
      +\textheight+\footskip\relax}{\lycee@licence}}}}}\fi}

\newcommand{\lycee@site}{faire-des-maths.fr}
\newcommand{\lycee@version}{{\scriptsize\color{gray}\lycee@site
  \ifdefined\lyceeversion\ --- Version~\lyceeversion\fi}}
\AtBeginDocument{\iflycee@licence\AddToHookNext{shipout/foreground}{%
  \put(0,0){\rlap{\kern\dimexpr1in+\hoffset+\oddsidemargin\relax
    \raisebox{-\dimexpr1in+\voffset+\topmargin+\headheight+\headsep
      +\textheight+\footskip\relax}{\lycee@version}}}}\fi}

\newcommand{\formulesserrees}{%
  \setlength{\abovedisplayskip}{3pt}\setlength{\belowdisplayskip}{3pt}%
  \setlength{\abovedisplayshortskip}{2pt}\setlength{\belowdisplayshortskip}{3pt}}

\setlength{\parindent}{0pt}

\RequirePackage[implicit=false,hidelinks,pdfencoding=auto,psdextra,bookmarksopen,
  bookmarksopenlevel=1,pdfcreator={},pdfproducer={}]{hyperref}
\RequirePackage{bookmark}
\hypersetup{pdfauthor={Luc Destombe},pdfkeywords={CC BY-NC-SA 4.0}}
\newcounter{lycee@signet}
\newcommand{\signet}[3][]{\stepcounter{lycee@signet}%
  \edef\lycee@dest{\if\relax\detokenize{#1}\relax
    signet.\arabic{lycee@signet}\else#1\fi}%
  \hypertarget{\lycee@dest}{}\bookmark[level=#2,dest=\lycee@dest]{#3}}
\pdfstringdefDisableCommands{%
  \def\R{R}\def\N{N}\def\Z{Z}\def\Q{Q}\def\vect#1{#1}\def\Cf{Cf}\def\Cg{Cg}\def\fois#1{×#1}%
  \def\barre#1{#1}\def\intervalle#1#2{[#1;#2]}\def\intff#1#2{[#1;#2]}%
  \def\textsuperscript#1{#1}\def\dfrac#1#2{#1/#2}\def\frac#1#2{#1/#2}%
  \def\vec#1{#1⃗}}%
\newcommand{\res}[1]{#1}
\geometry{top=2cm,bottom=1cm,left=1cm,right=1cm,headheight=14pt,footskip=0.6cm}

\RequirePackage{pgfplots}
\pgfplotsset{compat=1.18}
\RequirePackage{tkz-tab}

\newcounter{partiecours}
\renewcommand{\thepartiecours}{\Roman{partiecours}}
\newcommand{\partiecours}[1]{%
  \refstepcounter{partiecours}%
  \Needspace*{8\baselineskip}%
  \bigskip
  \noindent\begin{tcolorbox}[enhanced,colback=coulPart,colframe=coulPart,
      boxrule=0pt,arc=2pt,left=8pt,right=8pt,top=4pt,bottom=4pt,
      before skip=6pt,after skip=10pt]
    \color{white}\bfseries\large \thepartiecours{} --- #1%
    \signet[partie-\arabic{partiecours}]{1}{\thepartiecours{} --- #1}
  \end{tcolorbox}%
  \addcontentsline{toc}{section}{\thepartiecours{} --- #1}%
}

\newtcolorbox{boitecours}[3][]{%
  enhanced,breakable,
  colback=white,colframe=#2,coltitle=white,
  fonttitle=\bfseries,
  attach boxed title to top left={xshift=8pt,yshift=-\tcboxedtitleheight/2},
  boxed title style={colback=#2,arc=2pt,boxrule=0pt},
  boxrule=0.9pt,arc=2pt,
  left=8pt,right=8pt,top=8pt,bottom=6pt,
  title={#3},#1}

\newcounter{methode}
\newenvironment{definitionbox}[1][Définition]%
  {\begin{boitecours}[phantom={\signet{2}{#1}}]{coulDef}{#1}}{\end{boitecours}}
\newenvironment{proprietebox}[1][Propriété]%
  {\begin{boitecours}[phantom={\signet{2}{#1}}]{coulProp}{#1}}{\end{boitecours}}
\newenvironment{methodebox}[1][Méthode]{\stepcounter{methode}%
  \begin{boitecours}[phantom={\signet[methode-\arabic{methode}]{2}{#1}}]{coulMeth}{#1}}%
  {\end{boitecours}}

\newcommand{\etape}[1]{\par\smallskip\textbf{\color{coulMeth}#1}\par\nobreak\smallskip}
\newcommand{\enonce}[1]{\emph{#1}\par\smallskip}
\newcommand{\reponse}[1]{\par\smallskip
  {\footnotesize\itshape\color{black!55}Réponses : #1}}
\newcommand{\demo}[1]{\textbf{\color{coulMeth}Démonstration.} #1}
\newcommand{\afaire}[1]{%
  \par\smallskip
  \noindent{\small\color{coulExo}$\blacktriangleright$\ \textbf{À faire :}
    fiche d'exercices du chapitre, #1.}\par\smallskip}

\newenvironment{calculs}{\setlength{\jot}{5pt}\@fleqntrue\@mathmargin\retraitcalculs
  \setlength{\abovedisplayskip}{4pt}\setlength{\belowdisplayskip}{4pt}%
  \setlength{\abovedisplayshortskip}{4pt}\setlength{\belowdisplayshortskip}{4pt}%
  \csname alignat*\endcsname{2}}%
  {\csname endalignat*\endcsname}
\newcommand{\rem}[1]{\text{\small\itshape\color{black!60}#1}}
\newcommand{\explique}[2]{\par\smallskip\noindent
  \begin{minipage}[t]{0.57\linewidth}\raggedright #1\end{minipage}\hfill
  \begin{minipage}[t]{0.40\linewidth}\raggedright\leavevmode
    {\small\itshape\color{black!60}#2}\end{minipage}\par}
\newcommand{\cas}[1]{\par\smallskip\noindent\textbf{\color{coulExo}#1}\vadjust{\nobreak}}
\newcommand{\calc}[1]{\par\smallskip\textbf{\color{coulExo}#1}\par\nobreak}

\newtcolorbox{remarquebox}[1][Remarque]{%
  enhanced,breakable,colback=white,colframe=coulRem,
  boxrule=0pt,leftrule=2.5pt,arc=0pt,outer arc=0pt,
  left=8pt,right=6pt,top=5pt,bottom=5pt,
  before upper={\textbf{\color{coulRem}#1}\par\smallskip}}

\newtcolorbox{exemplebox}[1][Exemple]{%
  enhanced,breakable,colback=white,colframe=coulEx,
  boxrule=0pt,leftrule=2.5pt,arc=0pt,outer arc=0pt,
  left=8pt,right=6pt,top=5pt,bottom=5pt,
  before upper={\textbf{\color{coulEx}#1}\par\smallskip}}

\pgfplotsset{
  repere/.style={
    axis lines=middle,
    axis line style={-{Stealth[length=2.4mm]},thick},
    every axis x label/.style={at={(ticklabel* cs:1.0)},anchor=west},
    every axis y label/.style={at={(ticklabel* cs:1.0)},anchor=south},
    label style={font=\footnotesize},
    tick label style={font=\scriptsize},
    grid=both,
    major grid style={grille},
    minor tick num=0,
    tick align=inside,
    clip mode=individual,
  },
}
\tikzset{
  courbe/.style={coulCourbe,line width=1.1pt,smooth},
  courbeB/.style={coulCourbeB,line width=1.1pt,smooth},
  courbeC/.style={coulCourbeC,line width=1.1pt,smooth},
}

\newlist{etapes}{enumerate}{1}
\setlist[etapes]{label=\textbf{\arabic*.},leftmargin=*,itemsep=1pt,topsep=2pt}

\newlist{expressions}{enumerate}{1}
\setlist[expressions]{label={\Alph*\ $=$},leftmargin=*,itemsep=3pt,topsep=3pt}

\newcounter{exo}
\renewcommand{\theexo}{\ifnum\value{exo}<10 0\fi\arabic{exo}}
\newtcolorbox{exercice}[1][]{%
  enhanced,breakable,
  colback=white,colframe=coulExo,
  boxrule=0.9pt,arc=2pt,
  left=8pt,right=8pt,top=8pt,bottom=6pt,
  coltitle=white,fonttitle=\bfseries,
  attach boxed title to top left={xshift=8pt,yshift=-\tcboxedtitleheight/2},
  boxed title style={colback=coulExo,arc=2pt,boxrule=0pt},
  phantom={\signet{2}{Exercice \arabic{exo}}},
  title={Exercice \theexo\ifx\relax#1\relax\else{} --- #1\fi}}
\newenvironment{exo}[1][\relax]{\refstepcounter{exo}\begin{exercice}[#1]}{\end{exercice}}

  \newenvironment{correction}{\par}{\par}
  \newcommand{\autableau}{}
  \newcommand{\etapeexemples}{\etape{Exemples corrigés}}
  \newcommand{\etapetableau}[1]{\etape{#1}}
  \newenvironment{exempletableau}[1][Exemple]%
    {\begin{exemplebox}[#1]}{\end{exemplebox}}

\RequirePackage{comment}
\excludecomment{correctionavj}

\newcommand{\coursEntete}{}
\newcommand{\coursNiveau}{}
\newcommand{\coursTitre}{}
\newcommand{\coursSurTitre}{}
\newcommand{\coursSousTitre}{}

\newcommand{\enteteCours}[1]{\renewcommand{\coursEntete}{#1}}
\newcommand{\chapitrecours}[2]{\enteteCours{Chapitre #1 --- #2}}
\newcommand{\niveaucours}[1]{\renewcommand{\coursNiveau}{#1}}
\newcommand{\titrecours}[3][]{%
  \renewcommand{\coursTitre}{#2}%
  \renewcommand{\coursSurTitre}{#1}%
  \renewcommand{\coursSousTitre}{#3}}

\pagestyle{fancy}
\fancyhf{}
\fancyhead[L]{\footnotesize\color{coulPart}\textbf{\coursEntete}}
\fancyhead[R]{\footnotesize\color{coulPart}\textbf{\coursNiveau}}
\fancyfoot[C]{\footnotesize\thepage}
\renewcommand{\headrulewidth}{0.4pt}
\renewcommand{\headrule}{\hbox to\headwidth{\color{coulPart!40}\leaders\hrule height \headrulewidth\hfill}}

\setlength{\parskip}{2pt}

\newcommand{\lycee@mentioneleve}{\\[2pt]{\small\itshape Version élève}}
\AtBeginDocument{%
  \begin{center}
    \begin{tcolorbox}[enhanced,width=0.72\linewidth,colback=white,colframe=coulPart,
        boxrule=1.6pt,arc=3pt,halign=center,top=10pt,bottom=10pt]
      \ifx\coursSurTitre\empty
        {\Huge\bfseries\color{coulPart} \coursTitre}\\[4pt]
      \else
        {\Huge\bfseries\color{coulPart} \coursTitre}\\[3pt]
        {\Large\bfseries\color{coulPart} \coursSurTitre}\\[5pt]
      \fi
      {\normalsize \coursSousTitre}%
      
    \end{tcolorbox}
  \end{center}%
}

\RequirePackage{tikz}
\RequirePackage{xcolor}

\usetikzlibrary{decorations.text}
\definecolor{lycee@metal}{HTML}{1F4E79}
\definecolor{lycee@manche}{HTML}{8A5A2B}
\definecolor{lycee@etincelle}{HTML}{F28C28}
\newcommand{\lycee@pluriel}[2]{#1 #2\ifnum#1>1 s\fi}
\newcommand{\lycee@sceau}{%
  \edef\lycee@bas{\lycee@pluriel{\lyceerelectures}{relecture} et
    \lycee@pluriel{\lyceerodages}{rodage}}%
  \begin{tikzpicture}[line join=round]
    \fill[coulPart!10,draw=coulPart,line width=.8pt] (0:2.2)
      \foreach \i in {1,...,71}{ -- (\i*5:{isodd(\i)?2.1:2.2}) } -- cycle;
    \draw[coulPart,line width=1.2pt] (0,0) circle (1.93);
    \draw[coulPart,line width=.4pt] (0,0) circle (1.85);
    \draw[coulPart,line width=.4pt] (0,0) circle (1.2);
    \path[decorate,decoration={text along path,text align=center,
        text={|\fontsize{9.5}{11}\selectfont\bfseries\color{coulPart}|TEST{\'E} AVEC MES {\'E}L{\`E}VES}}]
      (180:1.45) arc (180:0:1.45);
    \path[decorate,decoration={text along path,text align=center,
        text/.expanded={|\noexpand\fontsize{8.5}{10}\noexpand\selectfont\noexpand\color{coulPart}|\lycee@bas}}]
      (195:1.68) arc (195:345:1.68);
    \begin{scope}[x=.25mm,y=-.25mm]
      \begin{scope}[shift={(-32,-40)}]%
        \fill[lycee@metal] (23,31) rectangle (41,49);
        \draw[lycee@etincelle,line width=.5pt,line cap=round]
          (20,30)--(15,27) (19,34)--(14,35) (44,30)--(49,27) (45,34)--(50,35);
        \fill[lycee@manche,rounded corners=.3mm] (38,24.6) rectangle (60,27.6);
        \fill[lycee@metal,rounded corners=.4mm] (25,20) rectangle (39,31);
        \fill[lycee@etincelle,rounded corners=.2mm] (25,20) rectangle (28,31);
      \end{scope}
    \end{scope}
    \ifdefined\lyceesceauplusieurs\def\lycee@annee{Ann{\'e}es~:}\else
      \def\lycee@annee{Ann{\'e}e}\fi
    \path[decorate,decoration={text along path,text align=center,
        text/.expanded={|\noexpand\fontsize{8}{10}\noexpand\selectfont\noexpand\color{coulPart}|%
          \unexpanded\expandafter{\lycee@annee} \unexpanded\expandafter{\lyceesceauannees}}}]
      (95:1.08) arc (95:445:1.08);
  \end{tikzpicture}}
\AtEndDocument{\ifdefined\lyceesceauannees
  \par\besoinplace{5cm}\vspace{1.5em}%
  \begin{center}\lycee@sceau\end{center}\fi}
\makeatother

\chapitrecours{8}{Trigonométrie}
\niveaucours{1\textsuperscript{re} spécialité mathématiques}
\titrecours{TRIGONOMÉTRIE}{Cours --- Première spécialité}

\begin{document}

\ShorthandsOff
\tikzset{
  lecture/.style={dashed,black!60,line width=0.6pt},
  sens/.style={-{Stealth[length=1.8mm]},line width=0.9pt,coulCourbeB},
  sensB/.style={-{Stealth[length=1.8mm]},line width=0.9pt,coulCourbeC},
}
\newcommand{\cerclenu}{%
  \draw[-{Stealth[length=2mm]},thick] (-1.3,0)--(1.3,0) node[below,font=\scriptsize,inner sep=2pt]{$x$};
  \draw[-{Stealth[length=2mm]},thick] (0,-1.3)--(0,1.3) node[right,font=\scriptsize,inner sep=2pt]{$y$};
  \draw[line width=0.9pt] (0,0) circle (1);
  \node[below left,font=\scriptsize,inner sep=1.5pt] at (0,0) {$O$};}
\newcommand{\cerclebase}{\cerclenu
  \node[below right,font=\scriptsize,inner sep=1.5pt] at (1,0) {$1$};
  \node[below left,font=\scriptsize,inner sep=1.5pt] at (-1,0) {$-1$};
  \node[above right,font=\scriptsize,inner sep=1.5pt] at (0,1) {$1$};
  \node[below right,font=\scriptsize,inner sep=1.5pt] at (0,-1) {$-1$};}
\newcommand{\cerclerep}{\cerclenu
  \node[above right,font=\scriptsize,inner sep=1.5pt] at (1,0) {$0$};
  \node[below left,font=\scriptsize,inner sep=1.5pt] at (1,0) {$2\pi$};
  \node[above left,font=\scriptsize,inner sep=1.5pt] at (-1,0) {$\pi$};}
\newcommand{\cerclearc}[3]{%
  \begin{tikzpicture}[x=0.95cm,y=0.95cm]
    \cerclenu
    \draw[-{Stealth[length=2.4mm]},red,line width=1.8pt] (1,0) arc (0:#1:1);
    \node[red,font=\footnotesize,inner sep=1pt,fill=white] at ({1.32*cos(#1+14)},{1.32*sin(#1+14)}) {#2};
  \end{tikzpicture}\par\smallskip{\small #3}}
\newcommand{\figasso}[5]{\begin{tikzpicture}[x=1.25cm,y=1.25cm]
  #5\cerclerep\ptc{#1}{#2}\ptc{#3}{#4}\end{tikzpicture}}
\newcommand{\rayons}[1]{\foreach \a in {#1}{%
  \draw[grille] (0,0)--({cos(\a)},{sin(\a)});}}
\newcommand{\ptc}[2]{%
  \path[coulCourbeB] ({cos(#1)},{sin(#1)}) node[croix]{};
  \node[font=\scriptsize,inner sep=1pt,align=center] at ({1.3*cos(#1)},{1.3*sin(#1)}) {#2};}
\newcommand{\verticale}[1]{\draw[lecture] ({cos(#1)},-1.15)--({cos(#1)},1.15);}
\newcommand{\horizontale}[1]{\draw[lecture] (-1.15,{sin(#1)})--(1.15,{sin(#1)});}
\newcommand{\lien}[2]{\draw[lecture] ({cos(#1)},{sin(#1)})--({cos(#2)},{sin(#2)});}
\newcommand{\lierx}[1]{\draw[lecture] ({cos(#1)},{sin(#1)})--({cos(#1)},0);}
\newcommand{\liery}[1]{\draw[lecture] ({cos(#1)},{sin(#1)})--(0,{sin(#1)});}
\newcommand{\valx}[2]{\pgfmathtruncatemacro{\q}{cos(#1)>0?1:0}%
  \draw[coulCourbeB,line width=1pt] ({cos(#1)},0.05)--({cos(#1)},-0.05);
  \ifnum\q=1 \node[below left,coulCourbeB,font=\scriptsize,inner sep=1pt,outer sep=1pt,fill=white] at ({cos(#1)},0) {#2};
  \else \node[below right,coulCourbeB,font=\scriptsize,inner sep=1pt,outer sep=1pt,fill=white] at ({cos(#1)},0) {#2};\fi}
\newcommand{\valy}[2]{%
  \draw[coulCourbeB,line width=1pt] (0.05,{sin(#1)})--(-0.05,{sin(#1)});
  \node[left,coulCourbeB,font=\scriptsize,inner sep=1pt,outer sep=1pt,fill=white] at (0,{sin(#1)}) {#2};}
\newcommand{\equacos}[4]{\begin{tikzpicture}[x=1.5cm,y=1.5cm]
  \cerclebase\verticale{#1}\valx{#1}{#4}\ptc{#1}{#2}\ptc{-#1}{#3}\end{tikzpicture}}
\newcommand{\equasin}[5]{\begin{tikzpicture}[x=1.5cm,y=1.5cm]
  \cerclebase\horizontale{#1}\valy{#1}{#5}\ptc{#1}{#3}\ptc{#2}{#4}\end{tikzpicture}}
\newcommand{\equapt}[4]{\begin{tikzpicture}[x=1.5cm,y=1.5cm]
  \cerclebase\verticale{#1}\horizontale{#1}\valx{#1}{#3}\valy{#1}{#4}\ptc{#1}{#2}\end{tikzpicture}}
\newcommand{\cerclept}[3]{\begin{tikzpicture}[x=1.5cm,y=1.5cm]
  \cerclebase#3\ptc{#1}{#2}\end{tikzpicture}}
\newcommand{\ptstable}[1]{\foreach \a in {-150,-135,-120,-90,-60,-45,-30,0,30,45,60,90,120,135,150,180}{%
  \path[coulCourbeB] ({\a*pi/180},{#1(\a)}) node[croix]{};}}
\newcommand{\arcsens}[4][sens]{%
  \draw[#1] ({#4*cos(#2)},{#4*sin(#2)}) arc (#2:#3:#4);}
\newcommand{\cerclemesure}[4]{%
  \begin{tikzpicture}[x=1.25cm,y=1.25cm]
    \rayons{#1}\cerclerep
    \arcsens{#3}{#2}{0.55}
    \ptc{#2}{#4}
  \end{tikzpicture}}
\newcommand{\figassociee}[3]{%
  \begin{tikzpicture}[x=1.25cm,y=1.25cm]
    #3\cerclerep
    \ptc{35}{$x$}
    \ptc{#1}{#2}
  \end{tikzpicture}}
\newcommand{\reperetrigo}[2]{%
  \foreach \k in {#1,...,#2}{\draw[grille] ({\k*pi/6},-1.5)--({\k*pi/6},1.5);}
  \foreach \j in {-3,...,3}{\draw[grille] ({#1*pi/6},{\j/2})--({#2*pi/6},{\j/2});}
  \draw[-{Stealth[length=2mm]},thick] ({#1*pi/6},0)--({#2*pi/6+0.4},0)
    node[below left=-1pt,font=\footnotesize]{$x$};
  \draw[-{Stealth[length=2mm]},thick] (0,-1.5)--(0,1.8)
    node[below left=-1pt,font=\footnotesize]{$y$};
  \node[above left,font=\scriptsize,inner sep=1pt,outer sep=1pt] at (0,1) {$1$};
  \foreach \j/\l in {0.5/\frac{1}{2},-0.5/-\frac{1}{2},-1/-1}{%
    \node[left,font=\scriptsize,inner sep=1pt,outer sep=1pt,fill=white] at (0,\j) {$\l$};}}
\newcommand{\gradpi}[2][below]{\foreach \k/\l in {#2}{%
  \draw ({\k*pi/6},0.05)--({\k*pi/6},-0.05);
  \node[#1,font=\scriptsize,inner sep=1pt,outer sep=2.5pt,fill=white] at ({\k*pi/6},0) {$\l$};}}
\newcommand{\gradcos}{%
  \gradpi{-12/-2\pi,-6/-\pi,6/\pi,12/2\pi}
  \gradpi[below left]{-9/-\frac{3\pi}{2},3/\frac{\pi}{2}}
  \gradpi[below right]{-3/-\frac{\pi}{2},9/\frac{3\pi}{2}}}
\newcommand{\gradsin}{%
  \gradpi{-9/-\frac{3\pi}{2},-3/-\frac{\pi}{2},3/\frac{\pi}{2},9/\frac{3\pi}{2}}
  \gradpi[below left]{-12/-2\pi,-6/-\pi,6/\pi}
  \gradpi[below right]{12/2\pi}}
\newcommand{\courbecos}[2]{\draw[coulCourbe,line width=1.1pt]
  plot[domain=#1*pi/6:#2*pi/6,samples=121] (\x,{cos(\x r)});}
\newcommand{\courbesin}[2]{\draw[coulCourbe,line width=1.1pt]
  plot[domain=#1*pi/6:#2*pi/6,samples=121] (\x,{sin(\x r)});}
\newcommand{\droiteh}[4]{\draw[coulCourbeB,line width=0.8pt]
  ({#1*pi/6},#3)--({#2*pi/6},#3) node[right,font=\footnotesize]{#4};
  \ifdim#3pt=0.5pt\node[left,font=\scriptsize,inner sep=1pt,outer sep=1pt,fill=white] at (0,0.5) {$\frac{1}{2}$};\fi
  \ifdim#3pt=-0.5pt\node[left,font=\scriptsize,inner sep=1pt,outer sep=1pt,fill=white] at (0,-0.5) {$-\frac{1}{2}$};\fi}
\newcommand{\lirex}[2]{\draw[lecture] ({#1*pi/6},#2)--({#1*pi/6},0);
  \path[coulCourbeB] ({#1*pi/6},#2) node[croix]{};}
\newcommand{\interaxe}[3][red]{\draw[#1,line width=1.8pt] ({#2*pi/6},0)--({#3*pi/6},0);}
\ShorthandsOn

\begin{remarquebox}[Comment utiliser ce cours]
Chaque partie commence par ce qu'il faut \textbf{savoir} (définitions, théorèmes), puis
vient ce qu'il faut \textbf{savoir faire} : une méthode, avec des
\textbf{exemples corrigés} faits en classe, puis des exercices
\og\textbf{À vous de jouer}\fg{} à chercher seul.

\end{remarquebox}

\partiecours{Cercle trigonométrique et radian}

\begin{definitionbox}[Définition 1 --- Cercle trigonométrique]
\begin{minipage}[c]{0.58\linewidth}\raggedright
Dans un repère orthonormé $(O\,;x\,;y)$, le \textbf{cercle trigonométrique} est le
cercle de centre $O$ et de rayon $1$, que l'on parcourt à partir du point de
coordonnées $(1\,;0)$.
\begin{itemize}
  \item Le \textbf{sens direct} (ou sens trigonométrique) est le sens contraire des
  aiguilles d'une montre.
  \item L'autre sens est le \textbf{sens indirect}.
\end{itemize}
\end{minipage}\hfill
\begin{minipage}[c]{0.40\linewidth}\centering
\begin{tikzpicture}[x=1.15cm,y=1.15cm]
  \cerclenu
  \draw[-{Stealth[length=2mm]},red,line width=1.1pt] (20:1.22) arc (20:70:1.22);
  \draw[-{Stealth[length=2mm]},line width=1.1pt] (-20:1.22) arc (-20:-70:1.22);
  \node[red,font=\footnotesize,anchor=south west,inner sep=1pt] at (48:1.32) {$+$ sens direct};
  \node[font=\footnotesize,anchor=north west,inner sep=1pt] at (-48:1.32) {$-$ sens indirect};
  \path[coulCourbeB] (1,0) node[croix]{};
  \draw[-{Stealth[length=1.6mm]},coulCourbeB,line width=0.8pt] (1.45,0.42)--(1.05,0.06);
  \node[coulCourbeB,font=\scriptsize,anchor=south west,inner sep=1pt] at (1.42,0.4) {départ};
\end{tikzpicture}
\end{minipage}
\end{definitionbox}

\begin{definitionbox}[Définition 2 --- Radian]
Un angle au centre mesure \textbf{un radian} (1 rad) quand il intercepte sur le cercle
trigonométrique un arc de longueur $1$. Plus généralement, la mesure en radians d'un
angle au centre est la \textbf{longueur de l'arc} qu'il intercepte sur ce cercle.

\smallskip
Le cercle trigonométrique a pour périmètre $2\pi\times 1$, soit $2\pi$ : un tour complet
mesure $360^{\circ}$ ou $2\pi$ rad. Un radian vaut à peu près $57^{\circ}$.

\medskip
\begin{minipage}[t]{0.24\linewidth}\centering
\cerclearc{90}{$\frac{\pi}{2}$}{un quart de tour}
\end{minipage}\hfill
\begin{minipage}[t]{0.24\linewidth}\centering
\cerclearc{180}{$\pi$}{un demi-tour}
\end{minipage}\hfill
\begin{minipage}[t]{0.24\linewidth}\centering
\cerclearc{270}{$\frac{3\pi}{2}$}{trois quarts de tour}
\end{minipage}\hfill
\begin{minipage}[t]{0.24\linewidth}\centering
\cerclearc{359.9}{$2\pi$}{un tour complet}
\end{minipage}
\end{definitionbox}

\begin{proprietebox}[Propriété 1 --- Degrés et radians]
\begin{minipage}[c]{0.50\linewidth}\raggedright
Les mesures d'un angle en degrés et en radians sont proportionnelles :
$180^{\circ}$ valent $\pi$ rad. Pour passer des degrés aux radians, on multiplie par
$\dfrac{\pi}{180}$ ; dans l'autre sens, par $\dfrac{180}{\pi}$.

\smallskip
Sur la figure : les degrés à l'extérieur du cercle, les radians à l'intérieur (en rouge).
\end{minipage}\hfill
\begin{minipage}[c]{0.48\linewidth}\centering
\begin{tikzpicture}[x=2.1cm,y=2.1cm]
  \foreach \a in {30,45,60}{\draw[grille] (0,0)--({cos(\a)},{sin(\a)});}
  \cerclenu
  \foreach \a/\d/\r in {30/30/\frac{\pi}{6},45/45/\frac{\pi}{4},60/60/\frac{\pi}{3},
      90/90/\frac{\pi}{2},180/180/\pi,270/270/\frac{3\pi}{2}}{%
    \path[coulCourbeB] ({cos(\a)},{sin(\a)}) node[croix]{};
    \node[font=\footnotesize,inner sep=1pt,fill=white] at ({1.17*cos(\a)},{1.17*sin(\a)}) {$\d^{\circ}$};
    \node[red,font=\footnotesize,inner sep=1pt,fill=white] at ({0.8*cos(\a)},{0.8*sin(\a)}) {$\r$};}
  \path[coulCourbeB] (1,0) node[croix]{};
  \node[font=\footnotesize,anchor=south west,inner sep=1pt] at (1.03,0.03) {$0^{\circ}$ ou $360^{\circ}$};
  \node[red,font=\footnotesize,anchor=south east,inner sep=1pt] at (0.97,0.03) {$0$ ou $2\pi$};
\end{tikzpicture}
\end{minipage}
\begin{center}
{\renewcommand{\arraystretch}{2.4}\setlength{\tabcolsep}{10pt}
\begin{tabular}{|m{1.5cm}|*{7}{>{\centering\arraybackslash}m{0.8cm}|}}
\hline
Degrés & $0$ & $30$ & $45$ & $60$ & $90$ & $180$ & $360$\\
\hline
Radians & $0$ & $\dfrac{\pi}{6}$ & $\dfrac{\pi}{4}$ & $\dfrac{\pi}{3}$ & $\dfrac{\pi}{2}$ & $\pi$ & $2\pi$\\
\hline
\end{tabular}}
\end{center}
\end{proprietebox}

\begin{methodebox}[Méthode 1 --- Convertir des degrés en radians et des radians en degrés]
\etapeexemples
Convertir chaque mesure d'angle :
\begin{questions}[label=\alph*)]
  \item $20^{\circ}$, en radians ;
  \item $135^{\circ}$, en radians ;
  \item $\dfrac{5\pi}{18}$ rad, en degrés ;
  \item $\dfrac{11\pi}{12}$ rad, en degrés.
\end{questions}
\begin{correction}
\cas{a)}
\begin{calculs}
20^{\circ} &= 20\times\dfrac{\pi}{180}\text{ rad} &\qquad& \rem{des degrés aux radians : on multiplie par $\tfrac{\pi}{180}$}\\
20^{\circ} &= \dfrac{\barre{20}\times\pi}{9\times\barre{20}}\text{ rad} && \rem{$180$ vaut $9\times 20$ : on simplifie par $20$}\\
20^{\circ} &= \res{\dfrac{\pi}{9}}\text{ rad} &&
\end{calculs}
\cas{b)}
\begin{calculs}
135^{\circ} &= 135\times\dfrac{\pi}{180}\text{ rad} &\qquad& \rem{même calcul}\\
135^{\circ} &= \dfrac{3\times\barre{45}\times\pi}{4\times\barre{45}}\text{ rad} && \rem{on simplifie par $45$}\\
135^{\circ} &= \res{\dfrac{3\pi}{4}}\text{ rad} &&
\end{calculs}
\cas{c)}
\begin{calculs}
\dfrac{5\pi}{18}\text{ rad} &= \dfrac{5\barre{\pi}}{18}\times\dfrac{180}{\barre{\pi}}\text{ degrés} &\qquad& \rem{des radians aux degrés : on multiplie par $\tfrac{180}{\pi}$}\\
\dfrac{5\pi}{18}\text{ rad} &= 5\times 10\text{ degrés} && \rem{$180$ divisé par $18$ vaut $10$}\\
\dfrac{5\pi}{18}\text{ rad} &= \res{50^{\circ}} &&
\end{calculs}
\cas{d)}
\begin{calculs}
\dfrac{11\pi}{12}\text{ rad} &= \dfrac{11\barre{\pi}}{12}\times\dfrac{180}{\barre{\pi}}\text{ degrés} &\qquad& \rem{même calcul}\\
\dfrac{11\pi}{12}\text{ rad} &= 11\times 15\text{ degrés} && \rem{$180$ divisé par $12$ vaut $15$}\\
\dfrac{11\pi}{12}\text{ rad} &= \res{165^{\circ}} &&
\end{calculs}
\end{correction}

\etape{À vous de jouer}
Convertir chaque mesure d'angle : a)~$80^{\circ}$, en radians ;
b)~$210^{\circ}$, en radians ; c)~$\dfrac{7\pi}{10}$ rad, en degrés ;
d)~$\dfrac{3\pi}{5}$ rad, en degrés.
\reponse{a)~$\dfrac{4\pi}{9}$ rad ;\; b)~$\dfrac{7\pi}{6}$ rad ;\; c)~$126^{\circ}$ ;\; d)~$108^{\circ}$.}
\end{methodebox}

\afaire{exercices 1 à 4}

\partiecours{Longueur d'un arc, mesures d'un angle orienté}

\begin{proprietebox}[Propriété 2 --- Longueur d'un arc]
\begin{minipage}[c]{0.62\linewidth}\raggedright
Sur un cercle de rayon $r$, un angle au centre de mesure $\alpha$ \textbf{radians}
intercepte un arc de longueur :
\[ \ell=r\times\alpha \]
La longueur $\ell$ s'exprime dans l'unité du rayon. Sur le cercle trigonométrique, de
rayon $1$, la longueur de l'arc est la mesure de l'angle en radians.
\end{minipage}\hfill
\begin{minipage}[c]{0.34\linewidth}\centering
\begin{tikzpicture}[x=1cm,y=1cm]
  \draw[line width=0.8pt] (0,0) circle (1.4);
  \draw (0,0)--(1.4,0);
  \draw (0,0)--({1.4*cos(65)},{1.4*sin(65)});
  \draw[red,line width=1.8pt] (1.4,0) arc (0:65:1.4);
  \draw (0.45,0) arc (0:65:0.45);
  \node[font=\footnotesize] at (32:0.72) {$\alpha$};
  \node[red,font=\footnotesize] at (32:1.75) {$\ell$};
  \node[below,font=\footnotesize] at (0.7,0) {$r$};
  \node[below left,font=\scriptsize,inner sep=1.5pt] at (0,0) {$O$};
\end{tikzpicture}
\end{minipage}
\end{proprietebox}

\begin{methodebox}[Méthode 2 --- Calculer la longueur d'un arc ou la mesure d'un angle]
\etapeexemples
\begin{questions}[label=\alph*)]
  \item L'aiguille des minutes d'une horloge mesure $15$ cm. En $5$ minutes, elle tourne
  de $\dfrac{\pi}{6}$ rad. Quelle longueur son extrémité parcourt-elle ?
  \item La corde d'un pendule mesure $80$ cm. La masse parcourt un arc de $20$ cm. De
  quel angle la corde a-t-elle tourné ?
\end{questions}
\begin{correction}
\cas{a)}
\begin{calculs}
\ell &= r\times\alpha &\qquad& \rem{l'angle est en radians}\\
\ell &= 15\times\dfrac{\pi}{6} && \rem{$r$ vaut $15$ cm et $\alpha$ vaut $\tfrac{\pi}{6}$}\\
\ell &= \dfrac{5\times\barre{3}\times\pi}{2\times\barre{3}} && \rem{on simplifie par $3$}\\
\ell &= \res{\dfrac{5\pi}{2}}\text{ cm} &&
\end{calculs}
\explique{L'extrémité parcourt environ $7{,}9$ cm.}{valeur approchée à la calculatrice}
\cas{b)}
\begin{calculs}
\ell &= r\times\alpha &\qquad& \rem{on cherche $\alpha$}\\
20 &= 80\times\alpha && \rem{$\ell$ vaut $20$ cm et $r$ vaut $80$ cm}\\
\alpha &= \dfrac{20}{80} && \rem{on divise par $80$}\\
\alpha &= \res{\dfrac{1}{4}} && \rem{on simplifie par $20$}
\end{calculs}
\explique{La corde a tourné de $\dfrac{1}{4}$ de radian.}{soit environ $14^{\circ}$}
\end{correction}

\etape{À vous de jouer}
\begin{questions}[label=\alph*)]
  \item Une pizza de rayon $12$ cm est coupée en $8$ parts égales, d'angle
  $\dfrac{\pi}{4}$ rad. Quelle est la longueur de croûte d'une part ?
  \item Un essuie-glace mesure $50$ cm. Son extrémité parcourt un arc de $20$ cm. De
  quel angle a-t-il tourné ?
\end{questions}
\reponse{a)~$3\pi$ cm, soit environ $9{,}4$ cm ;\; b)~$\dfrac{2}{5}$ rad.}
\end{methodebox}

\begin{definitionbox}[Définition 3 --- Mesure d'un angle orienté, image d'un nombre réel]
\begin{minipage}[c]{0.66\linewidth}\raggedright
On part du point de coordonnées $(1\,;0)$ et on tourne jusqu'au point $M$ du cercle
trigonométrique. La longueur parcourue est une \textbf{mesure de l'angle orienté} formé
par le demi-axe $[Ox)$ et $[OM)$ :
\begin{itemize}
  \item positive dans le sens direct ;
  \item négative dans le sens indirect.
\end{itemize}
Si $x$ est cette mesure, on dit que $M$ est l'\textbf{image} du nombre réel $x$ : c'est
comme si l'on enroulait la droite des nombres réels autour du cercle.

\smallskip
Chaque point du cercle a une mesure positive et une mesure négative, qui aboutissent au
même endroit : par exemple, $\dfrac{\pi}{2}$ et $-\dfrac{3\pi}{2}$ pour le point $M$.
\end{minipage}\hfill
\begin{minipage}[c]{0.30\linewidth}\centering
\begin{tikzpicture}[x=1.25cm,y=1.25cm]
  \cerclerep
  \arcsens{0}{90}{0.72}
  \arcsens[sensB]{0}{-270}{0.45}
  \ptc{90}{$M\quad\ $}
  \node[coulCourbeB,font=\scriptsize] at (0.78,0.78) {$\frac{\pi}{2}$};
  \node[coulCourbeC,font=\scriptsize] at (-0.55,-0.55) {$-\frac{3\pi}{2}$};
\end{tikzpicture}
\end{minipage}
\end{definitionbox}

\begin{proprietebox}[Propriété 3 --- Plusieurs mesures pour un même point]
Faire un tour de plus ou de moins ramène au même point : deux nombres réels $x$ et $x'$
ont la même image si et seulement si
\[ x'=x+2k\pi,\quad\text{avec }k\in\mathbb{Z}\quad(k\text{ est le nombre de tours}). \]
\end{proprietebox}

\begin{methodebox}[Méthode 3 --- Lire les mesures d'un point du cercle]
\etapeexemples
\begin{minipage}[c]{0.40\linewidth}\raggedright
Pour chaque point $A$ et $B$ de la figure, lire :
\begin{questions}[label=\alph*)]
  \item sa mesure dans le sens direct, entre $0$ et $2\pi$ ;
  \item sa mesure dans le sens indirect, entre $-2\pi$ et $0$.
\end{questions}
\end{minipage}\hfill
\begin{minipage}[c]{0.28\linewidth}\centering
Douzièmes de tour\par
\begin{tikzpicture}[x=1.25cm,y=1.25cm]
  \rayons{0,30,...,330}\cerclerep
  \ptc{150}{$A$}
\end{tikzpicture}
\end{minipage}\hfill
\begin{minipage}[c]{0.28\linewidth}\centering
Huitièmes de tour\par
\begin{tikzpicture}[x=1.25cm,y=1.25cm]
  \rayons{0,45,...,315}\cerclerep
  \ptc{315}{$B$}
\end{tikzpicture}
\end{minipage}
\begin{correction}
\cas{a)}
\explique{Pour $A$ : $5\times\dfrac{\pi}{6}$, soit $\res{\dfrac{5\pi}{6}}$.}{un douzième de tour vaut $\tfrac{\pi}{6}$ ; on en compte $5$ à partir du point $(1\,;0)$}
\explique{Pour $B$ : $7\times\dfrac{\pi}{4}$, soit $\res{\dfrac{7\pi}{4}}$.}{un huitième de tour vaut $\tfrac{\pi}{4}$ ; on en compte $7$}
\cas{b)}
\explique{Pour $A$ : $7\times\left(-\dfrac{\pi}{6}\right)$, soit $\res{-\dfrac{7\pi}{6}}$.}{dans l'autre sens, on compte $7$ douzièmes de tour}
\explique{Pour $B$ : $\res{-\dfrac{\pi}{4}}$.}{dans l'autre sens, un seul huitième de tour}
\end{correction}

\etape{À vous de jouer}
\begin{minipage}[c]{0.40\linewidth}\raggedright
Pour chaque point $C$ et $D$ de la figure, lire :
\begin{questions}[label=\alph*)]
  \item sa mesure dans le sens direct, entre $0$ et $2\pi$ ;
  \item sa mesure dans le sens indirect, entre $-2\pi$ et $0$.
\end{questions}
\end{minipage}\hfill
\begin{minipage}[c]{0.28\linewidth}\centering
Sixièmes de tour\par
\begin{tikzpicture}[x=1.25cm,y=1.25cm]
  \rayons{0,60,...,300}\cerclerep
  \ptc{240}{$C$}
\end{tikzpicture}
\end{minipage}\hfill
\begin{minipage}[c]{0.28\linewidth}\centering
Huitièmes de tour\par
\begin{tikzpicture}[x=1.25cm,y=1.25cm]
  \rayons{0,45,...,315}\cerclerep
  \ptc{135}{$D$}
\end{tikzpicture}
\end{minipage}
\reponse{a)~$C$ : $\dfrac{4\pi}{3}$ ; $D$ : $\dfrac{3\pi}{4}$ ;\; b)~$C$ : $-\dfrac{2\pi}{3}$ ;
$D$ : $-\dfrac{5\pi}{4}$.}
\end{methodebox}

\begin{methodebox}[Méthode 4 --- Placer sur le cercle l'image d'un nombre réel]
\etapeexemples
Placer sur le cercle trigonométrique l'image de chaque nombre réel :
\begin{center}
$A$, image de $\dfrac{3\pi}{4}$ \qquad $B$, image de $-\dfrac{5\pi}{6}$ \qquad
$C$, image de $\dfrac{7\pi}{3}$ \qquad $D$, image de $-\dfrac{9\pi}{2}$
\end{center}
\begin{correction}
\cas{a)}\par\nopagebreak
\begin{minipage}[c]{0.68\linewidth}\raggedright
\explique{On compte $3$ fois $\dfrac{\pi}{4}$ à partir du point $(1\,;0)$.}{positif : sens direct}
\end{minipage}\hfill
\begin{minipage}[c]{0.28\linewidth}\centering
\cerclemesure{0,45,...,315}{135}{0}{$A$}
\end{minipage}
\cas{b)}\par\nopagebreak
\begin{minipage}[c]{0.68\linewidth}\raggedright
\explique{On compte $5$ fois $\dfrac{\pi}{6}$ à partir du point $(1\,;0)$.}{négatif : sens indirect}
\end{minipage}\hfill
\begin{minipage}[c]{0.28\linewidth}\centering
\cerclemesure{0,30,...,330}{-150}{0}{$B$}
\end{minipage}
\cas{c)}\par\nopagebreak
\begin{minipage}[c]{0.68\linewidth}\raggedright
\begin{calculs}
\dfrac{7\pi}{3} &= \dfrac{6\pi}{3}+\dfrac{\pi}{3} &\quad& \rem{on fait apparaître des tours}\\
\dfrac{7\pi}{3} &= 2\pi+\dfrac{\pi}{3} && \rem{$\tfrac{6\pi}{3}$ vaut $2\pi$ : un tour}
\end{calculs}
\explique{$C$ est l'image de $\res{\dfrac{\pi}{3}}$.}{propriété 3, avec $k=1$}
\end{minipage}\hfill
\begin{minipage}[c]{0.28\linewidth}\centering
\cerclemesure{0,60,...,300}{60}{0}{$C$}
\end{minipage}
\cas{d)}\par\nopagebreak
\begin{minipage}[c]{0.68\linewidth}\raggedright
\begin{calculs}
-\dfrac{9\pi}{2} &= -\dfrac{8\pi}{2}-\dfrac{\pi}{2} &\quad& \rem{$-9$ s'écrit $-8-1$}\\
-\dfrac{9\pi}{2} &= -4\pi-\dfrac{\pi}{2} && \rem{$-4\pi$ : deux tours en sens indirect}
\end{calculs}
\explique{$D$ est l'image de $\res{-\dfrac{\pi}{2}}$.}{propriété 3, avec $k=-2$}
\end{minipage}\hfill
\begin{minipage}[c]{0.28\linewidth}\centering
\cerclemesure{0,90,180,270}{-90}{0}{$D$}
\end{minipage}
\end{correction}

\etape{À vous de jouer}
Placer sur le cercle trigonométrique l'image de chaque nombre réel :
\begin{center}
$E$, image de $\dfrac{5\pi}{4}$ \qquad $F$, image de $-\dfrac{\pi}{3}$ \qquad
$G$, image de $\dfrac{19\pi}{6}$ \qquad $H$, image de $-\dfrac{7\pi}{2}$
\end{center}
\reponse{$G$ est l'image de $\dfrac{7\pi}{6}$ ;\; $H$ est l'image de $\dfrac{\pi}{2}$.}
\end{methodebox}

\afaire{exercices 5 à 9}

\partiecours{Mesure principale d'un angle orienté}

Un même point a une infinité de mesures : on en choisit une, toujours la même.

\begin{definitionbox}[Définition 4 --- Mesure principale]
\begin{minipage}[c]{0.66\linewidth}\raggedright
La \textbf{mesure principale} d'un angle orienté est sa seule mesure qui appartient à
l'intervalle $\intof{-\pi}{\pi}$.

\smallskip
Les points du haut du cercle ont une mesure principale positive, ceux du bas une mesure
principale négative.
\end{minipage}\hfill
\begin{minipage}[c]{0.30\linewidth}\centering
\begin{tikzpicture}[x=1.25cm,y=1.25cm]
  \cerclerep
  \draw[coulCourbeB,line width=1.5pt] (1,0) arc (0:180:1);
  \draw[coulCourbeC,line width=1.5pt] (1,0) arc (0:-180:1);
  \node[coulCourbeB,font=\scriptsize] at (0.45,0.45) {$+$};
  \node[coulCourbeC,font=\scriptsize] at (0.45,-0.45) {$-$};
\end{tikzpicture}
\end{minipage}
\end{definitionbox}

\begin{methodebox}[Méthode 5 --- Déterminer la mesure principale d'un angle orienté]
\etapeexemples
Déterminer la mesure principale de chaque angle orienté, de mesure :
\begin{center}
a)~$\dfrac{11\pi}{4}$ \qquad b)~$-\dfrac{9\pi}{5}$ \qquad c)~$\dfrac{43\pi}{6}$ \qquad d)~$-15\pi$
\end{center}
\begin{correction}
\cas{a)}
\begin{calculs}
2\pi &= \dfrac{8\pi}{4} &\qquad& \rem{un tour, écrit avec le dénominateur $4$}\\
\dfrac{11\pi}{4}-\dfrac{8\pi}{4} &= \dfrac{3\pi}{4} && \rem{on retire un tour : même point du cercle}
\end{calculs}
\explique{La mesure principale est $\res{\dfrac{3\pi}{4}}$.}{$\tfrac{3\pi}{4}$ appartient bien à $\intof{-\pi}{\pi}$}
\cas{b)}
\begin{calculs}
2\pi &= \dfrac{10\pi}{5} &\qquad& \rem{un tour, écrit avec le dénominateur $5$}\\
-\dfrac{9\pi}{5}+\dfrac{10\pi}{5} &= \dfrac{\pi}{5} && \rem{mesure négative : on ajoute un tour}
\end{calculs}
\explique{La mesure principale est $\res{\dfrac{\pi}{5}}$.}{$\tfrac{\pi}{5}$ appartient bien à $\intof{-\pi}{\pi}$}
\cas{c)}
\begin{calculs}
2\pi &= \dfrac{12\pi}{6} &\qquad& \rem{un tour, écrit avec le dénominateur $6$}\\
4\times 2\pi &= \dfrac{48\pi}{6} && \rem{$48$ : le multiple de $12$ le plus proche de $43$}\\
\dfrac{43\pi}{6}-\dfrac{48\pi}{6} &= -\dfrac{5\pi}{6} && \rem{on retire $4$ tours}
\end{calculs}
\explique{La mesure principale est $\res{-\dfrac{5\pi}{6}}$.}{$-\tfrac{5\pi}{6}$ appartient à $\intof{-\pi}{\pi}$ ; avec $36$, on trouverait $\tfrac{7\pi}{6}$, trop grand}
\cas{d)}
\begin{calculs}
-15\pi+16\pi &= \pi &\qquad& \rem{$16\pi$ : $8$ tours}
\end{calculs}
\explique{La mesure principale est $\res{\pi}$.}{$\pi$ appartient à $\intof{-\pi}{\pi}$, mais pas $-\pi$}
\end{correction}

\etape{À vous de jouer}
Déterminer la mesure principale de chaque angle orienté, de mesure :
\begin{center}
a)~$\dfrac{17\pi}{3}$ \qquad b)~$-\dfrac{13\pi}{4}$ \qquad c)~$\dfrac{101\pi}{6}$ \qquad d)~$-11\pi$
\end{center}
\reponse{a)~$-\dfrac{\pi}{3}$ ;\; b)~$\dfrac{3\pi}{4}$ ;\; c)~$\dfrac{5\pi}{6}$ ;\; d)~$\pi$.}
\end{methodebox}

\afaire{exercices 10 à 13}

\partiecours{Cosinus et sinus d'un nombre réel}

\begin{definitionbox}[Définition 5 --- Cosinus et sinus d'un nombre réel]
\begin{minipage}[c]{0.66\linewidth}\raggedright
Soit $M$ l'image du nombre réel $x$ sur le cercle trigonométrique.
\begin{itemize}
  \item Le \textbf{cosinus} de $x$, noté $\cos(x)$, est l'abscisse de $M$.
  \item Le \textbf{sinus} de $x$, noté $\sin(x)$, est l'ordonnée de $M$.
\end{itemize}
Entre $0$ et $\dfrac{\pi}{2}$, on retrouve le cosinus et le sinus du triangle rectangle
$OHM$, d'angle $\widehat{HOM}$ de mesure $x$, car l'hypoténuse $OM$ vaut $1$ :
\begin{calculs}
\cos\big(\widehat{HOM}\big) &= \dfrac{OH}{OM} &\qquad& \rem{côté adjacent sur hypoténuse}\\
\cos\big(\widehat{HOM}\big) &= OH && \rem{$OM$ vaut $1$}
\end{calculs}
De même, $\sin\big(\widehat{HOM}\big)=\dfrac{HM}{OM}$, soit $\sin\big(\widehat{HOM}\big)=HM$.
\end{minipage}\hfill
\begin{minipage}[c]{0.30\linewidth}\centering
\begin{tikzpicture}[x=1.7cm,y=1.7cm]
  \cerclebase
  \draw[red,line width=0.8pt] (0.3,0) arc (0:50:0.3);
  \node[red,font=\scriptsize] at (25:0.42) {$x$};
  \draw[black!70] (0,0)--({cos(50)},{sin(50)});
  \draw[lecture] ({cos(50)},{sin(50)})--({cos(50)},0);
  \draw[lecture] ({cos(50)},{sin(50)})--(0,{sin(50)});
  \draw[coulCourbeB,line width=1.6pt] (0,0)--({cos(50)},0);
  \draw[coulCourbeC,line width=1.6pt] (0,0)--(0,{sin(50)});
  \ptc{50}{$M$}
  \node[font=\scriptsize,above right,inner sep=1pt] at ({cos(50)},0) {$H$};
  \node[coulCourbeB,font=\scriptsize,below] at ({cos(50)/2},0) {$\cos(x)$};
  \node[coulCourbeC,font=\scriptsize,left] at (0,{sin(50)/2}) {$\sin(x)$};
\end{tikzpicture}
\end{minipage}
\end{definitionbox}

\begin{proprietebox}[Propriété 4 --- Encadrement, relation, tours]
Pour tout nombre réel $x$ et tout entier relatif $k$ :
\begin{center}
$-1\le\cos(x)\le 1$ \qquad $-1\le\sin(x)\le 1$ \qquad $\cos^{2}(x)+\sin^{2}(x)=1$\\[4pt]
$\cos(x+2k\pi)=\cos(x)$ \qquad $\sin(x+2k\pi)=\sin(x)$
\end{center}
où $\cos^{2}(x)$ désigne $\big(\cos(x)\big)^{2}$.
\demo{le point $M$ est sur un cercle de rayon $1$ : ses coordonnées sont entre $-1$ et $1$.
Les nombres réels $x$ et $x+2k\pi$ ont la même image (propriété 3), donc le même cosinus
et le même sinus. Enfin, le triangle $OHM$ est rectangle en $H$ :}
\begin{calculs}
OH^{2}+HM^{2} &= OM^{2} &\qquad& \rem{théorème de Pythagore}\\
\cos^{2}(x)+\sin^{2}(x) &= 1 && \rem{$OM$ vaut $1$}
\end{calculs}
\end{proprietebox}

\begin{proprietebox}[Propriété 5 --- Valeurs remarquables]
\begin{minipage}[c]{0.60\linewidth}\centering
{\renewcommand{\arraystretch}{2.4}\setlength{\tabcolsep}{7pt}
\begin{tabular}{|>{\centering\arraybackslash}m{1.2cm}|*{6}{>{\centering\arraybackslash}m{0.65cm}|}}
\hline
$x$ & $0$ & $\dfrac{\pi}{6}$ & $\dfrac{\pi}{4}$ & $\dfrac{\pi}{3}$ & $\dfrac{\pi}{2}$ & $\pi$\\
\hline
$\cos(x)$ & $1$ & $\dfrac{\sqrt{3}}{2}$ & $\dfrac{\sqrt{2}}{2}$ & $\dfrac{1}{2}$ & $0$ & $-1$\\
\hline
$\sin(x)$ & $0$ & $\dfrac{1}{2}$ & $\dfrac{\sqrt{2}}{2}$ & $\dfrac{\sqrt{3}}{2}$ & $1$ & $0$\\
\hline
\end{tabular}}

\medskip
{\raggedright À retenir avec l'image du quart de cercle : le cosinus diminue de $1$ à $0$
quand $x$ augmente, le sinus augmente de $0$ à $1$.

\smallskip
Pour situer les valeurs : $\dfrac{\sqrt{2}}{2}\approx 0{,}71$ et
$\dfrac{\sqrt{3}}{2}\approx 0{,}87$, donc $\dfrac{1}{2}<\dfrac{\sqrt{2}}{2}<\dfrac{\sqrt{3}}{2}$.\par}
\end{minipage}\hfill
\begin{minipage}[c]{0.36\linewidth}\centering
\begin{tikzpicture}[x=4cm,y=4cm]
  \draw[-{Stealth[length=2mm]},thick] (0,0)--(1.2,0);
  \draw[-{Stealth[length=2mm]},thick] (0,0)--(0,1.2);
  \draw[line width=0.9pt] (1,0) arc (0:90:1);
  \foreach \a/\l/\c/\s in {30/\frac{\pi}{6}/\frac{\sqrt{3}}{2}/\frac{1}{2},
      45/\frac{\pi}{4}/\frac{\sqrt{2}}{2}/\frac{\sqrt{2}}{2},
      60/\frac{\pi}{3}/\frac{1}{2}/\frac{\sqrt{3}}{2}}{%
    \draw[lecture] ({cos(\a)},0)--({cos(\a)},{sin(\a)})--(0,{sin(\a)});
    \draw[grille] (0,0)--({cos(\a)},{sin(\a)});
    \path[coulCourbeB] ({cos(\a)},{sin(\a)}) node[croix]{};
    \node[font=\scriptsize] at ({1.12*cos(\a)},{1.12*sin(\a)}) {$\l$};
    \node[coulCourbeB,font=\scriptsize,below,inner sep=2pt] at ({cos(\a)},0) {$\c$};
    \node[coulCourbeC,font=\scriptsize,left,inner sep=2pt] at (0,{sin(\a)}) {$\s$};}
  \node[font=\scriptsize,below,inner sep=2pt] at (1,0) {$1$};
  \node[font=\scriptsize,left,inner sep=2pt] at (0,1) {$1$};
  \node[font=\scriptsize,below left,inner sep=2pt] at (0,0) {$O$};
\end{tikzpicture}
\end{minipage}

\medskip
Les valeurs en $0$, $\dfrac{\pi}{2}$ et $\pi$ se lisent sur le cercle ; celles en $\dfrac{\pi}{6}$, $\dfrac{\pi}{4}$ et $\dfrac{\pi}{3}$ se démontrent
dans la fiche d'exercices.
\end{proprietebox}

\begin{methodebox}[Méthode 6 --- Déterminer un cosinus ou un sinus à l'aide de la mesure principale]
\etapeexemples
Déterminer la valeur exacte de :
\begin{center}
a)~$\sin\left(\dfrac{\pi}{6}\right)$ \qquad b)~$\cos\left(\dfrac{25\pi}{4}\right)$ \qquad
c)~$\sin\left(-\dfrac{11\pi}{6}\right)$ \qquad d)~$\cos(-9\pi)$
\end{center}
\begin{correction}
\cas{a)}
\explique{$\sin\left(\dfrac{\pi}{6}\right)=\res{\dfrac{1}{2}}$.}{lecture du tableau}
\cas{b)}
\begin{calculs}
\dfrac{25\pi}{4}-\dfrac{24\pi}{4} &= \dfrac{\pi}{4} &\qquad& \rem{on retire $3$ tours : même point}\\
\cos\left(\dfrac{25\pi}{4}\right) &= \cos\left(\dfrac{\pi}{4}\right) && \rem{même point, même cosinus}\\
\cos\left(\dfrac{25\pi}{4}\right) &= \res{\dfrac{\sqrt{2}}{2}} && \rem{lecture du tableau}
\end{calculs}
\cas{c)}
\begin{calculs}
-\dfrac{11\pi}{6}+\dfrac{12\pi}{6} &= \dfrac{\pi}{6} &\qquad& \rem{on ajoute un tour}\\
\sin\left(-\dfrac{11\pi}{6}\right) &= \sin\left(\dfrac{\pi}{6}\right) && \rem{même point, même sinus}\\
\sin\left(-\dfrac{11\pi}{6}\right) &= \res{\dfrac{1}{2}} && \rem{lecture du tableau}
\end{calculs}
\cas{d)}
\begin{calculs}
-9\pi+10\pi &= \pi &\qquad& \rem{on ajoute $5$ tours}\\
\cos(-9\pi) &= \cos(\pi) && \rem{même point, même cosinus}\\
\cos(-9\pi) &= \res{-1} && \rem{lecture du tableau}
\end{calculs}
\end{correction}

\etape{À vous de jouer}
Déterminer la valeur exacte de :
\begin{center}
a)~$\cos\left(\dfrac{\pi}{6}\right)$ \qquad b)~$\sin\left(\dfrac{19\pi}{3}\right)$ \qquad
c)~$\cos\left(-\dfrac{7\pi}{4}\right)$ \qquad d)~$\cos(-7\pi)$
\end{center}
\reponse{a)~$\dfrac{\sqrt{3}}{2}$ ;\; b)~$\dfrac{\sqrt{3}}{2}$ ;\; c)~$\dfrac{\sqrt{2}}{2}$ ;\; d)~$-1$.}
\end{methodebox}

\begin{methodebox}[Méthode 7 --- Calculer une expression avec des cosinus et des sinus]
\etapeexemples
Calculer la valeur exacte de :
\begin{questions}[label=\alph*)]
  \item $f(x)=2\cos(x)+4\sin(x)$, pour $x=\dfrac{\pi}{3}$ ;
  \item $g(x)=4\cos^{2}(x)-1$, pour $x=\dfrac{\pi}{6}$ ;
  \item $h(x)=\cos(3x)$, pour $x=\dfrac{\pi}{9}$.
\end{questions}
\begin{correction}
\cas{a)}
\begin{calculs}
f\left(\dfrac{\pi}{3}\right) &= 2\cos\left(\dfrac{\pi}{3}\right)+4\sin\left(\dfrac{\pi}{3}\right) &\qquad& \rem{on remplace $x$ par $\tfrac{\pi}{3}$}\\
f\left(\dfrac{\pi}{3}\right) &= 2\times\dfrac{1}{2}+4\times\dfrac{\sqrt{3}}{2} && \rem{valeurs du tableau}\\
f\left(\dfrac{\pi}{3}\right) &= \res{1+2\sqrt{3}} && \rem{$4$ divisé par $2$ vaut $2$}
\end{calculs}
\cas{b)}
\begin{calculs}
g\left(\dfrac{\pi}{6}\right) &= 4\left(\cos\left(\dfrac{\pi}{6}\right)\right)^{2}-1 &\qquad& \rem{$\cos^{2}$ : le cosinus, au carré}\\
g\left(\dfrac{\pi}{6}\right) &= 4\times\left(\dfrac{\sqrt{3}}{2}\right)^{2}-1 && \rem{valeur du tableau}\\
g\left(\dfrac{\pi}{6}\right) &= 4\times\dfrac{3}{4}-1 && \rem{$\big(\sqrt{3}\big)^{2}$ vaut $3$ et $2^{2}$ vaut $4$}\\
g\left(\dfrac{\pi}{6}\right) &= \res{2} && \rem{$3-1$}
\end{calculs}
\cas{c)}
\begin{calculs}
h\left(\dfrac{\pi}{9}\right) &= \cos\left(3\times\dfrac{\pi}{9}\right) &\qquad& \rem{on multiplie d'abord, on prend le cosinus ensuite}\\
h\left(\dfrac{\pi}{9}\right) &= \cos\left(\dfrac{\pi}{3}\right) && \rem{$3\times\tfrac{\pi}{9}$ vaut $\tfrac{\pi}{3}$}\\
h\left(\dfrac{\pi}{9}\right) &= \res{\dfrac{1}{2}} && \rem{valeur du tableau}
\end{calculs}
\end{correction}

\etape{À vous de jouer}
Calculer la valeur exacte de :
\begin{questions}[label=\alph*)]
  \item $k(x)=6\cos(x)-2\sin(x)$, pour $x=\dfrac{\pi}{6}$ ;
  \item $m(x)=1-2\sin^{2}(x)$, pour $x=\dfrac{\pi}{4}$ ;
  \item $n(x)=\sin(4x)$, pour $x=\dfrac{\pi}{8}$.
\end{questions}
\reponse{a)~$3\sqrt{3}-1$ ;\; b)~$0$ ;\; c)~$1$.}
\end{methodebox}

\afaire{exercices 14 à 18}

\partiecours{Angles associés}

Le tableau s'arrête à $\pi$ et ne donne rien pour $\dfrac{2\pi}{3}$ ou $-\dfrac{\pi}{3}$ :
ces points sont les symétriques de points connus, leurs coordonnées s'en déduisent.

\begin{proprietebox}[Propriété 6 --- Angles associés]
Pour tout nombre réel $x$, on lit sur le cercle :

\medskip
\begin{minipage}[t]{0.32\linewidth}\centering
\figassociee{-35}{$-x$}{\draw[coulCourbeC,line width=0.8pt] (-1.2,0)--(1.2,0);\lien{35}{-35}\liery{35}\liery{-35}}\par
{\small symétrie d'axe $(Ox)$ :\\ même cosinus, sinus opposés}\par\smallskip
$\cos(-x)=\cos(x)$\par
$\sin(-x)=-\sin(x)$
\end{minipage}\hfill
\begin{minipage}[t]{0.32\linewidth}\centering
\figassociee{145}{$\pi-x$}{\draw[coulCourbeC,line width=0.8pt] (0,-1.2)--(0,1.2);\lien{35}{145}\lierx{35}\lierx{145}}\par
{\small symétrie d'axe $(Oy)$ :\\ même sinus, cosinus opposés}\par\smallskip
$\cos(\pi-x)=-\cos(x)$\par
$\sin(\pi-x)=\sin(x)$
\end{minipage}\hfill
\begin{minipage}[t]{0.32\linewidth}\centering
\figassociee{215}{$\pi+x$}{\lien{35}{215}\lierx{35}\liery{35}\lierx{215}\liery{215}}\par
{\small symétrie de centre $O$ :\\ cosinus opposés, sinus opposés}\par\smallskip
$\cos(\pi+x)=-\cos(x)$\par
$\sin(\pi+x)=-\sin(x)$
\end{minipage}

\medskip
\bigskip
\hfill
\begin{minipage}[t]{0.32\linewidth}\centering
\figassociee{55}{$\frac{\pi}{2}-x\quad$}{\draw[coulCourbeC,line width=0.8pt] (-1.1,-1.1)--(1.1,1.1);\lien{35}{55}\liery{35}\lierx{55}}\par
{\small symétrie d'axe la droite $y=x$ :\\ cosinus et sinus s'échangent}\par\smallskip
$\cos\left(\dfrac{\pi}{2}-x\right)=\sin(x)$\par\smallskip
$\sin\left(\dfrac{\pi}{2}-x\right)=\cos(x)$
\end{minipage}\hspace{1.5cm}
\begin{minipage}[t]{0.32\linewidth}\centering
\figassociee{125}{$\frac{\pi}{2}+x$}{\draw[lecture] (0,0)--({cos(35)},{sin(35)}); \draw[lecture] (0,0)--({cos(125)},{sin(125)});\liery{35}\lierx{125}}\par
{\small quart de tour de centre $O$ :\\ ils s'échangent, un signe change}\par\smallskip
$\cos\left(\dfrac{\pi}{2}+x\right)=-\sin(x)$\par\smallskip
$\sin\left(\dfrac{\pi}{2}+x\right)=\cos(x)$
\end{minipage}
\hfill\null

\medskip
Ces formules ne s'apprennent pas une à une : on dessine la figure et on regarde si
l'abscisse et l'ordonnée changent de signe ou s'échangent.

\end{proprietebox}

\begin{methodebox}[Méthode 8 --- Calculer un cosinus ou un sinus à l'aide d'un angle associé]
\etapeexemples
Déterminer la valeur exacte de :
\begin{center}
a)~$\cos\left(\dfrac{2\pi}{3}\right)$ \qquad b)~$\sin\left(\dfrac{5\pi}{6}\right)$ \qquad
c)~$\cos\left(\dfrac{5\pi}{4}\right)$ \qquad d)~$\sin\left(-\dfrac{\pi}{3}\right)$
\end{center}
\begin{correction}
\cas{a)}\par\nopagebreak
\begin{minipage}[c]{0.68\linewidth}\raggedright
\begin{calculs}
\cos\left(\dfrac{2\pi}{3}\right) &= \cos\left(\pi-\dfrac{\pi}{3}\right) &\qquad& \rem{$\tfrac{3\pi}{3}-\tfrac{\pi}{3}$ vaut $\tfrac{2\pi}{3}$}\\
\cos\left(\dfrac{2\pi}{3}\right) &= -\cos\left(\dfrac{\pi}{3}\right) && \rem{$\cos(\pi-x)=-\cos(x)$}\\
\cos\left(\dfrac{2\pi}{3}\right) &= \res{-\dfrac{1}{2}} && \rem{valeur du tableau}
\end{calculs}
\end{minipage}\hfill
\begin{minipage}[c]{0.28\linewidth}\centering
\figasso{60}{$\frac{\pi}{3}$}{120}{$\frac{2\pi}{3}$}{\lien{60}{120}\lierx{60}\lierx{120}}
\end{minipage}
\cas{b)}\par\nopagebreak
\begin{minipage}[c]{0.68\linewidth}\raggedright
\begin{calculs}
\sin\left(\dfrac{5\pi}{6}\right) &= \sin\left(\pi-\dfrac{\pi}{6}\right) &\qquad& \rem{$\tfrac{6\pi}{6}-\tfrac{\pi}{6}$ vaut $\tfrac{5\pi}{6}$}\\
\sin\left(\dfrac{5\pi}{6}\right) &= \sin\left(\dfrac{\pi}{6}\right) && \rem{$\sin(\pi-x)=\sin(x)$}\\
\sin\left(\dfrac{5\pi}{6}\right) &= \res{\dfrac{1}{2}} && \rem{valeur du tableau}
\end{calculs}
\end{minipage}\hfill
\begin{minipage}[c]{0.28\linewidth}\centering
\figasso{30}{$\frac{\pi}{6}$}{150}{$\frac{5\pi}{6}$}{\lien{30}{150}\liery{30}}
\end{minipage}
\cas{c)}\par\nopagebreak
\begin{minipage}[c]{0.68\linewidth}\raggedright
\begin{calculs}
\cos\left(\dfrac{5\pi}{4}\right) &= \cos\left(\pi+\dfrac{\pi}{4}\right) &\qquad& \rem{$\tfrac{4\pi}{4}+\tfrac{\pi}{4}$ vaut $\tfrac{5\pi}{4}$}\\
\cos\left(\dfrac{5\pi}{4}\right) &= -\cos\left(\dfrac{\pi}{4}\right) && \rem{$\cos(\pi+x)=-\cos(x)$}\\
\cos\left(\dfrac{5\pi}{4}\right) &= \res{-\dfrac{\sqrt{2}}{2}} && \rem{valeur du tableau}
\end{calculs}
\end{minipage}\hfill
\begin{minipage}[c]{0.28\linewidth}\centering
\figasso{45}{$\frac{\pi}{4}$}{225}{$\frac{5\pi}{4}$}{\lien{45}{225}\lierx{45}\lierx{225}}
\end{minipage}
\cas{d)}\par\nopagebreak
\begin{minipage}[c]{0.68\linewidth}\raggedright
\begin{calculs}
\sin\left(-\dfrac{\pi}{3}\right) &= -\sin\left(\dfrac{\pi}{3}\right) &\qquad& \rem{$\sin(-x)=-\sin(x)$}\\
\sin\left(-\dfrac{\pi}{3}\right) &= \res{-\dfrac{\sqrt{3}}{2}} && \rem{valeur du tableau}
\end{calculs}
\end{minipage}\hfill
\begin{minipage}[c]{0.28\linewidth}\centering
\figasso{60}{$\frac{\pi}{3}$}{-60}{$-\frac{\pi}{3}$}{\lien{60}{-60}\liery{60}\liery{-60}}
\end{minipage}
\end{correction}

\etape{À vous de jouer}
Déterminer la valeur exacte de :
\begin{center}
a)~$\cos\left(\dfrac{3\pi}{4}\right)$ \qquad b)~$\sin\left(\dfrac{2\pi}{3}\right)$ \qquad
c)~$\cos\left(\dfrac{7\pi}{6}\right)$ \qquad d)~$\sin\left(-\dfrac{\pi}{4}\right)$
\end{center}
\reponse{a)~$-\dfrac{\sqrt{2}}{2}$ ;\; b)~$\dfrac{\sqrt{3}}{2}$ ;\; c)~$-\dfrac{\sqrt{3}}{2}$ ;\; d)~$-\dfrac{\sqrt{2}}{2}$.}
\end{methodebox}

\afaire{exercices 19 à 22}

\partiecours{Simplifier une expression à l'aide des angles associés}

Les formules de la propriété 6, et celles des tours de la propriété 4, servent aussi
quand $x$ n'a pas de valeur : on remplace chaque terme, puis on réduit.

\begin{methodebox}[Méthode 9 --- Simplifier une expression à l'aide des angles associés]
\etapeexemples
Pour tout nombre réel $x$, simplifier :
\begin{center}
$A(x)=\sin(\pi-x)+\sin(-x)$ \qquad
$B(x)=\cos\left(\dfrac{\pi}{2}-x\right)+\sin(\pi+x)+\cos(-x)$\\[4pt]
$C(x)=2\sin\left(\dfrac{\pi}{2}+x\right)-\cos(\pi-x)+\sin(x+2\pi)$
\end{center}
\begin{correction}
\cas{Expression $A$}
\begin{calculs}
A(x) &= \sin(\pi-x)+\sin(-x) &\qquad& \rem{on recopie}\\
A(x) &= \sin(x)-\sin(x) && \rem{on remplace chaque terme, d'après la propriété 6}\\
A(x) &= \res{0} &&
\end{calculs}
\cas{Expression $B$}
\begin{calculs}
B(x) &= \cos\left(\dfrac{\pi}{2}-x\right)+\sin(\pi+x)+\cos(-x) &\qquad& \rem{on recopie}\\
B(x) &= \sin(x)-\sin(x)+\cos(x) && \rem{on remplace chaque terme}\\
B(x) &= \res{\cos(x)} && \rem{les sinus s'annulent}
\end{calculs}
\cas{Expression $C$}
\begin{calculs}
C(x) &= 2\sin\left(\dfrac{\pi}{2}+x\right)-\cos(\pi-x)+\sin(x+2\pi) &\qquad& \rem{on recopie}\\
C(x) &= 2\cos(x)-\big(-\cos(x)\big)+\sin(x) && \rem{$\sin(x+2\pi)=\sin(x)$ : un tour de plus}\\
C(x) &= \res{3\cos(x)+\sin(x)} && \rem{$2+1$ vaut $3$}
\end{calculs}
\end{correction}

\etape{À vous de jouer}
Pour tout nombre réel $x$, simplifier :
\begin{center}
$D(x)=\cos(\pi+x)+\cos(\pi-x)$ \qquad
$E(x)=\sin\left(\dfrac{\pi}{2}-x\right)+\cos(\pi+x)+\sin(-x)$\\[4pt]
$F(x)=3\cos\left(\dfrac{\pi}{2}+x\right)-\sin(\pi-x)+\cos(x+2\pi)$
\end{center}
\reponse{$D(x)=-2\cos(x)$ ;\; $E(x)=-\sin(x)$ ;\; $F(x)=-4\sin(x)+\cos(x)$.}
\end{methodebox}

\afaire{exercices 23 à 26}

\partiecours{Le cercle trigonométrique complet}

\begin{proprietebox}[Le cercle complet et son tableau de valeurs]
Les symétries donnent seize points, chacun avec sa mesure principale. Les pointillés
verticaux relient les points de même cosinus, les pointillés horizontaux ceux de même sinus.
Ce cercle est l'outil de tout le chapitre : \textbf{il faut savoir le refaire}.
\begin{center}
\begin{tikzpicture}[x=3.7cm,y=3.7cm]
  \foreach \a in {30,45,60}{%
    \draw[lecture] ({cos(\a)},{-sin(\a)})--({cos(\a)},{sin(\a)});
    \draw[lecture] ({-cos(\a)},{-sin(\a)})--({-cos(\a)},{sin(\a)});
    \draw[lecture] ({-cos(\a)},{sin(\a)})--({cos(\a)},{sin(\a)});
    \draw[lecture] ({-cos(\a)},{-sin(\a)})--({cos(\a)},{-sin(\a)});}
  \cerclebase
  \path[coulCourbeB] (1,0) node[croix]{} (0,1) node[croix]{} (-1,0) node[croix]{} (0,-1) node[croix]{};
  \node[font=\footnotesize,inner sep=1pt] at (1.1,0.09) {$0$};
  \node[font=\footnotesize,inner sep=1pt] at (-1.1,0.09) {$\pi$};
  \node[font=\footnotesize,inner sep=1pt] at (-0.12,1.12) {$\frac{\pi}{2}$};
  \node[font=\footnotesize,inner sep=1pt] at (-0.16,-1.12) {$-\frac{\pi}{2}$};
  \foreach \a/\l in {30/\frac{\pi}{6},45/\frac{\pi}{4},60/\frac{\pi}{3},
     120/\frac{2\pi}{3},135/\frac{3\pi}{4},150/\frac{5\pi}{6},
     -30/-\frac{\pi}{6},-45/-\frac{\pi}{4},-60/-\frac{\pi}{3},
     -120/-\frac{2\pi}{3},-135/-\frac{3\pi}{4},-150/-\frac{5\pi}{6}}{%
    \path[coulCourbeB] ({cos(\a)},{sin(\a)}) node[croix]{};
    \node[font=\footnotesize,inner sep=1pt] at ({1.13*cos(\a)},{1.13*sin(\a)}) {$\l$};}
  \foreach \v/\l in {0.5/\frac{1}{2},0.866/\frac{\sqrt{3}}{2},-0.5/-\frac{1}{2},-0.866/-\frac{\sqrt{3}}{2}}{%
    \draw (\v,0.02)--(\v,-0.02);
    \node[below,font=\footnotesize,inner sep=1pt,outer sep=2pt,coulCourbeB,fill=white] at (\v,0) {$\l$};
    \draw (0.02,\v)--(-0.02,\v);
    \node[left,font=\footnotesize,inner sep=1pt,outer sep=2pt,coulCourbeC,fill=white] at (0,\v) {$\l$};}
  \foreach \v/\l in {0.707/\frac{\sqrt{2}}{2},-0.707/-\frac{\sqrt{2}}{2}}{%
    \draw (\v,0.02)--(\v,-0.02);
    \node[above,font=\footnotesize,inner sep=1pt,outer sep=2pt,coulCourbeB,fill=white] at (\v,0) {$\l$};
    \draw (0.02,\v)--(-0.02,\v);
    \node[right,font=\footnotesize,inner sep=1pt,outer sep=2pt,coulCourbeC,fill=white] at (0,\v) {$\l$};}
\end{tikzpicture}
\end{center}
\begin{center}
{\renewcommand{\arraystretch}{2.4}\setlength{\tabcolsep}{4pt}
\begin{tabular}{|>{\centering\arraybackslash}m{1.2cm}|*{9}{>{\centering\arraybackslash}m{1.15cm}|}}
\hline
$x$ & $0$ & $\dfrac{\pi}{6}$ & $\dfrac{\pi}{4}$ & $\dfrac{\pi}{3}$ & $\dfrac{\pi}{2}$ & $\dfrac{2\pi}{3}$ & $\dfrac{3\pi}{4}$ & $\dfrac{5\pi}{6}$ & $\pi$\\
\hline
$\cos(x)$ & $1$ & $\dfrac{\sqrt{3}}{2}$ & $\dfrac{\sqrt{2}}{2}$ & $\dfrac{1}{2}$ & $0$ & $-\dfrac{1}{2}$ & $-\dfrac{\sqrt{2}}{2}$ & $-\dfrac{\sqrt{3}}{2}$ & $-1$\\
\hline
$\sin(x)$ & $0$ & $\dfrac{1}{2}$ & $\dfrac{\sqrt{2}}{2}$ & $\dfrac{\sqrt{3}}{2}$ & $1$ & $\dfrac{\sqrt{3}}{2}$ & $\dfrac{\sqrt{2}}{2}$ & $\dfrac{1}{2}$ & $0$\\
\hline
\end{tabular}

\medskip
\renewcommand{\arraystretch}{2.4}
\begin{tabular}{|>{\centering\arraybackslash}m{1.2cm}|*{7}{>{\centering\arraybackslash}m{1.15cm}|}}
\hline
$x$ & $-\dfrac{\pi}{6}$ & $-\dfrac{\pi}{4}$ & $-\dfrac{\pi}{3}$ & $-\dfrac{\pi}{2}$ & $-\dfrac{2\pi}{3}$ & $-\dfrac{3\pi}{4}$ & $-\dfrac{5\pi}{6}$\\
\hline
$\cos(x)$ & $\dfrac{\sqrt{3}}{2}$ & $\dfrac{\sqrt{2}}{2}$ & $\dfrac{1}{2}$ & $0$ & $-\dfrac{1}{2}$ & $-\dfrac{\sqrt{2}}{2}$ & $-\dfrac{\sqrt{3}}{2}$\\
\hline
$\sin(x)$ & $-\dfrac{1}{2}$ & $-\dfrac{\sqrt{2}}{2}$ & $-\dfrac{\sqrt{3}}{2}$ & $-1$ & $-\dfrac{\sqrt{3}}{2}$ & $-\dfrac{\sqrt{2}}{2}$ & $-\dfrac{1}{2}$\\
\hline
\end{tabular}}
\end{center}
\end{proprietebox}

\begin{methodebox}[Méthode 10 --- Retrouver la mesure principale d'après le cosinus et le sinus]
\etapeexemples
Déterminer la mesure principale de l'angle orienté de mesure $x$ tel que :
\begin{questions}[label=\alph*)]
  \item $\cos(x)=-\dfrac{\sqrt{2}}{2}$ et $\sin(x)=\dfrac{\sqrt{2}}{2}$ ;
  \item $\cos(x)=\dfrac{\sqrt{3}}{2}$ et $\sin(x)=-\dfrac{1}{2}$ ;
  \item $\cos(x)=-\dfrac{1}{2}$ et $\sin(x)=-\dfrac{\sqrt{3}}{2}$.
\end{questions}
\begin{correction}
\cas{a)}\par\nopagebreak
\begin{minipage}[c]{0.68\linewidth}\raggedright
\explique{On place $-\dfrac{\sqrt{2}}{2}$ sur l'axe des abscisses et $\dfrac{\sqrt{2}}{2}$ sur l'axe des ordonnées.}{on trace les deux droites}
\explique{Elles se croisent sur le cercle, au point associé à $\dfrac{3\pi}{4}$.}{un seul point a ces deux coordonnées}
\explique{$x=\res{\dfrac{3\pi}{4}}$.}{lecture sur le cercle complet}
\end{minipage}\hfill
\begin{minipage}[c]{0.28\linewidth}\centering
\equapt{135}{$\frac{3\pi}{4}$}{$-\frac{\sqrt{2}}{2}$}{$\frac{\sqrt{2}}{2}$}
\end{minipage}
\cas{b)}\par\nopagebreak
\begin{minipage}[c]{0.68\linewidth}\raggedright
\explique{On place $\dfrac{\sqrt{3}}{2}$ sur l'axe des abscisses et $-\dfrac{1}{2}$ sur l'axe des ordonnées.}{on trace les deux droites}
\explique{Elles se croisent sur le cercle, au point associé à $-\dfrac{\pi}{6}$.}{un seul point a ces deux coordonnées}
\explique{$x=\res{-\dfrac{\pi}{6}}$.}{lecture sur le cercle complet}
\end{minipage}\hfill
\begin{minipage}[c]{0.28\linewidth}\centering
\equapt{-30}{$-\frac{\pi}{6}$}{$\frac{\sqrt{3}}{2}$}{$-\frac{1}{2}$}
\end{minipage}
\cas{c)}\par\nopagebreak
\begin{minipage}[c]{0.68\linewidth}\raggedright
\explique{On place $-\dfrac{1}{2}$ sur l'axe des abscisses et $-\dfrac{\sqrt{3}}{2}$ sur l'axe des ordonnées.}{on trace les deux droites}
\explique{Elles se croisent sur le cercle, au point associé à $-\dfrac{2\pi}{3}$.}{en bas à gauche du cercle}
\explique{$x=\res{-\dfrac{2\pi}{3}}$.}{lecture sur le cercle complet}
\end{minipage}\hfill
\begin{minipage}[c]{0.28\linewidth}\centering
\equapt{-120}{$-\frac{2\pi}{3}$}{$-\frac{1}{2}$}{$-\frac{\sqrt{3}}{2}$}
\end{minipage}
\end{correction}

\etape{À vous de jouer}
Déterminer la mesure principale de l'angle orienté de mesure $x$ tel que :
\begin{questions}[label=\alph*)]
  \item $\cos(x)=-\dfrac{\sqrt{3}}{2}$ et $\sin(x)=-\dfrac{1}{2}$ ;
  \item $\cos(x)=\dfrac{1}{2}$ et $\sin(x)=\dfrac{\sqrt{3}}{2}$ ;
  \item $\cos(x)=-\dfrac{\sqrt{2}}{2}$ et $\sin(x)=-\dfrac{\sqrt{2}}{2}$.
\end{questions}
\reponse{a)~$-\dfrac{5\pi}{6}$ ;\; b)~$\dfrac{\pi}{3}$ ;\; c)~$-\dfrac{3\pi}{4}$.}
\end{methodebox}

\afaire{exercices 27 à 30}

\partiecours{Équations avec des cosinus et des sinus}

Résoudre $\cos(x)=a$, c'est chercher les points du cercle d'abscisse $a$, sur une droite
verticale ; résoudre $\sin(x)=a$, c'est chercher les points d'ordonnée $a$, sur une droite
horizontale. On travaille avec le \textbf{cercle complet} de la partie VII sous les yeux.

\begin{methodebox}[Méthode 11 --- Résoudre $\cos(x)=a$ ou $\sin(x)=a$ sur un intervalle]
\etapeexemples
Résoudre, à l'aide du cercle complet, sur l'intervalle $\intof{-\pi}{\pi}$ :
\begin{center}
a)~$\cos(x)=\dfrac{\sqrt{3}}{2}$ \qquad b)~$\sin(x)=\dfrac{1}{2}$ \qquad
c)~$\cos(x)=-\dfrac{\sqrt{2}}{2}$ \qquad d)~$\sin(x)=-1$ \qquad e)~$\cos(x)=\dfrac{3}{2}$
\end{center}
puis sur un autre intervalle :
\begin{center}
f)~$\cos(x)=-\dfrac{1}{2}$ sur $\intff{0}{2\pi}$ \qquad
g)~$\sin(x)=-\dfrac{\sqrt{2}}{2}$ sur $\intff{0}{2\pi}$\\[4pt]
h)~$\sin(x)=\dfrac{\sqrt{3}}{2}$ sur $\intff{-2\pi}{2\pi}$
\end{center}
\begin{correction}
\cas{a)}\par\nopagebreak
\begin{minipage}[c]{0.68\linewidth}\raggedright
\explique{On place $\dfrac{\sqrt{3}}{2}$ sur l'axe des abscisses et on trace la droite verticale.}{le cosinus se lit sur l'axe des abscisses}
\explique{Elle coupe le cercle en deux points, associés à $\dfrac{\pi}{6}$ et $-\dfrac{\pi}{6}$.}{lecture sur le cercle complet}
\explique{$S=\res{\ens{-\dfrac{\pi}{6}\,;\dfrac{\pi}{6}}}$.}{une solution par point}
\end{minipage}\hfill
\begin{minipage}[c]{0.28\linewidth}\centering
\equacos{30}{$\frac{\pi}{6}$}{$-\frac{\pi}{6}$}{$\frac{\sqrt{3}}{2}$}
\end{minipage}
\cas{b)}\par\nopagebreak
\begin{minipage}[c]{0.68\linewidth}\raggedright
\explique{On place $\dfrac{1}{2}$ sur l'axe des ordonnées et on trace la droite horizontale.}{le sinus se lit sur l'axe des ordonnées}
\explique{Elle coupe le cercle en deux points, associés à $\dfrac{\pi}{6}$ et $\dfrac{5\pi}{6}$.}{lecture sur le cercle complet}
\explique{$S=\res{\ens{\dfrac{\pi}{6}\,;\dfrac{5\pi}{6}}}$.}{une solution par point}
\end{minipage}\hfill
\begin{minipage}[c]{0.28\linewidth}\centering
\equasin{30}{150}{$\frac{\pi}{6}$}{$\frac{5\pi}{6}$}{$\frac{1}{2}$}
\end{minipage}
\cas{c)}\par\nopagebreak
\begin{minipage}[c]{0.68\linewidth}\raggedright
\explique{On place $-\dfrac{\sqrt{2}}{2}$ sur l'axe des abscisses et on trace la droite verticale.}{valeur négative : à gauche de $O$}
\explique{Elle coupe le cercle en deux points, associés à $\dfrac{3\pi}{4}$ et $-\dfrac{3\pi}{4}$.}{lecture sur le cercle complet}
\explique{$S=\res{\ens{-\dfrac{3\pi}{4}\,;\dfrac{3\pi}{4}}}$.}{une solution par point}
\end{minipage}\hfill
\begin{minipage}[c]{0.28\linewidth}\centering
\equacos{135}{$\frac{3\pi}{4}$}{$-\frac{3\pi}{4}$}{$-\frac{\sqrt{2}}{2}$}
\end{minipage}
\cas{d)}\par\nopagebreak
\begin{minipage}[c]{0.68\linewidth}\raggedright
\explique{La droite horizontale passant par $-1$ touche le cercle en un seul point.}{le point le plus bas du cercle}
\explique{$S=\res{\ens{-\dfrac{\pi}{2}}}$.}{une seule solution}
\end{minipage}\hfill
\begin{minipage}[c]{0.28\linewidth}\centering
\cerclept{-90}{$-\frac{\pi}{2}\qquad$}{\horizontale{-90}}
\end{minipage}
\cas{e)}
\explique{$S=\res{\varnothing}$ : il n'y a aucune solution.}{$\tfrac{3}{2}$ vaut $1{,}5$ ; la droite verticale passant par $1{,}5$ ne rencontre pas le cercle : un cosinus ne dépasse jamais $1$}
\cas{f)}\par\nopagebreak
\begin{minipage}[c]{0.68\linewidth}\raggedright
\explique{On place $-\dfrac{1}{2}$ sur l'axe des abscisses et on trace la droite verticale.}{le cosinus se lit sur l'axe des abscisses}
\explique{Elle coupe le cercle en deux points, de mesures principales $\dfrac{2\pi}{3}$ et $-\dfrac{2\pi}{3}$.}{lecture sur le cercle complet}
\explique{$\dfrac{2\pi}{3}$ appartient à $\intff{0}{2\pi}$, mais pas $-\dfrac{2\pi}{3}$.}{on compare à l'intervalle demandé}
\begin{calculs}
-\dfrac{2\pi}{3}+\dfrac{6\pi}{3} &= \dfrac{4\pi}{3} &\quad& \rem{on ajoute un tour : même point}
\end{calculs}
\explique{$S=\res{\ens{\dfrac{2\pi}{3}\,;\dfrac{4\pi}{3}}}$.}{une solution par point, dans l'intervalle demandé}
\end{minipage}\hfill
\begin{minipage}[c]{0.28\linewidth}\centering
\equacos{120}{$\frac{2\pi}{3}$}{$\frac{4\pi}{3}$}{$-\frac{1}{2}$}
\end{minipage}
\cas{g)}\par\nopagebreak
\begin{minipage}[c]{0.68\linewidth}\raggedright
\explique{On place $-\dfrac{\sqrt{2}}{2}$ sur l'axe des ordonnées et on trace la droite horizontale.}{le sinus se lit sur l'axe des ordonnées}
\explique{Elle coupe le cercle en deux points, de mesures principales $-\dfrac{\pi}{4}$ et $-\dfrac{3\pi}{4}$.}{lecture sur le cercle complet}
\explique{Aucune n'appartient à $\intff{0}{2\pi}$.}{elles sont négatives : on ajoute un tour à chacune}
\begin{calculs}
-\dfrac{\pi}{4}+\dfrac{8\pi}{4} &= \dfrac{7\pi}{4} &\quad& \rem{même point, autre mesure}\\
-\dfrac{3\pi}{4}+\dfrac{8\pi}{4} &= \dfrac{5\pi}{4} && \rem{même point, autre mesure}
\end{calculs}
\explique{$S=\res{\ens{\dfrac{5\pi}{4}\,;\dfrac{7\pi}{4}}}$.}{on range les solutions}
\end{minipage}\hfill
\begin{minipage}[c]{0.28\linewidth}\centering
\equasin{-45}{-135}{$\frac{7\pi}{4}$}{$\frac{5\pi}{4}$}{$-\frac{\sqrt{2}}{2}$}
\end{minipage}
\cas{h)}\par\nopagebreak
\begin{minipage}[c]{0.68\linewidth}\raggedright
\explique{On place $\dfrac{\sqrt{3}}{2}$ sur l'axe des ordonnées et on trace la droite horizontale.}{le sinus se lit sur l'axe des ordonnées}
\explique{Elle coupe le cercle en deux points, de mesures principales $\dfrac{\pi}{3}$ et $\dfrac{2\pi}{3}$.}{lecture sur le cercle complet}
\explique{L'intervalle $\intff{-2\pi}{2\pi}$ couvre deux tours : chaque point a une seconde mesure.}{deux tours : deux solutions par point}
\begin{calculs}
\dfrac{\pi}{3}-\dfrac{6\pi}{3} &= -\dfrac{5\pi}{3} &\quad& \rem{un tour de moins}\\
\dfrac{2\pi}{3}-\dfrac{6\pi}{3} &= -\dfrac{4\pi}{3} && \rem{un tour de moins}
\end{calculs}
\explique{$S=\res{\ens{-\dfrac{5\pi}{3}\,;-\dfrac{4\pi}{3}\,;\dfrac{\pi}{3}\,;\dfrac{2\pi}{3}}}$.}{quatre solutions}
\end{minipage}\hfill
\begin{minipage}[c]{0.28\linewidth}\centering
\equasin{60}{120}{$\frac{\pi}{3}$\\$-\frac{5\pi}{3}$}{$\frac{2\pi}{3}$\\$-\frac{4\pi}{3}$}{$\frac{\sqrt{3}}{2}$}
\end{minipage}
\end{correction}

\etape{À vous de jouer}
Résoudre, à l'aide du cercle complet, sur l'intervalle $\intof{-\pi}{\pi}$ :
\begin{center}
a)~$\cos(x)=-\dfrac{\sqrt{3}}{2}$ \qquad b)~$\sin(x)=\dfrac{\sqrt{2}}{2}$ \qquad
c)~$\cos(x)=\dfrac{1}{2}$ \qquad d)~$\sin(x)=1$ \qquad e)~$\sin(x)=-2$
\end{center}
puis sur un autre intervalle :
\begin{center}
f)~$\cos(x)=\dfrac{\sqrt{2}}{2}$ sur $\intff{0}{2\pi}$ \qquad
g)~$\sin(x)=-\dfrac{1}{2}$ sur $\intff{0}{2\pi}$\\[4pt]
h)~$\sin(x)=\dfrac{1}{2}$ sur $\intff{-2\pi}{2\pi}$
\end{center}
\reponse{a)~$\ens{-\dfrac{5\pi}{6}\,;\dfrac{5\pi}{6}}$ ;\; b)~$\ens{\dfrac{\pi}{4}\,;\dfrac{3\pi}{4}}$ ;\;
c)~$\ens{-\dfrac{\pi}{3}\,;\dfrac{\pi}{3}}$ ;\; d)~$\ens{\dfrac{\pi}{2}}$ ;\; e)~aucune solution ;\;
f)~$\ens{\dfrac{\pi}{4}\,;\dfrac{7\pi}{4}}$ ;\; g)~$\ens{\dfrac{7\pi}{6}\,;\dfrac{11\pi}{6}}$ ;\;
h)~$\ens{-\dfrac{11\pi}{6}\,;-\dfrac{7\pi}{6}\,;\dfrac{\pi}{6}\,;\dfrac{5\pi}{6}}$.}
\end{methodebox}

\afaire{exercices 31 à 34}

\partiecours{Fonctions cosinus et sinus}

\begin{definitionbox}[Définition 6 --- Fonctions cosinus et sinus, période]
La \textbf{fonction cosinus} est la fonction définie sur $\R$ par $x\longmapsto\cos(x)$ ;
la \textbf{fonction sinus}, la fonction définie sur $\R$ par $x\longmapsto\sin(x)$. Leurs
courbes s'appellent des \textbf{sinusoïdes}.

\smallskip
Soit $T$ un nombre réel strictement positif. Une fonction $f$ définie sur $\R$ est
\textbf{périodique de période $T$} si, pour tout nombre réel $x$ :
\[ f(x+T)=f(x) \]
Sa courbe reproduit alors le même motif sur chaque intervalle de longueur $T$.
\end{definitionbox}

\begin{proprietebox}[Propriété 7 --- Période et variations]
Chaque colonne du tableau de valeurs de la partie VII donne un point de la courbe (les
croix) : par exemple, $\cos\left(\dfrac{\pi}{3}\right)=\dfrac{1}{2}$ donne le point de
coordonnées $\left(\dfrac{\pi}{3}\,;\dfrac{1}{2}\right)$ sur la courbe du cosinus.

\smallskip
Comme $\cos(x+2\pi)=\cos(x)$ et $\sin(x+2\pi)=\sin(x)$ pour tout nombre réel $x$
(propriété 4), les fonctions cosinus et sinus sont \textbf{périodiques de période $2\pi$} :
le motif de $\intff{-\pi}{\pi}$ se répète. Dans les repères ci-dessous, un carreau vaut
$\dfrac{\pi}{6}$ en abscisse et $\dfrac{1}{2}$ en ordonnée.
\begin{center}
\begin{tikzpicture}[x=1.2cm,y=1.2cm]
  \reperetrigo{-12}{12}\courbecos{-12}{12}\gradcos\ptstable{cos}
  \node[coulCourbe,font=\footnotesize,fill=white,inner sep=1pt] at (2.2,1.25) {$y=\cos(x)$};
\end{tikzpicture}

\medskip
\begin{tikzpicture}[x=1.2cm,y=1.2cm]
  \reperetrigo{-12}{12}\courbesin{-12}{12}\gradsin\ptstable{sin}
  \node[coulCourbe,font=\footnotesize,fill=white,inner sep=1pt] at (3.3,1.25) {$y=\sin(x)$};
\end{tikzpicture}
\end{center}
Le tableau de valeurs donne aussi le sens de variation : quand $x$ va de $-\pi$ à $0$,
$\cos(x)$ monte de $-1$ à $1$, puis redescend à $-1$ quand $x$ va de $0$ à $\pi$.
Tableaux de variation sur une période, $\intff{-\pi}{\pi}$ :

\medskip
\begin{minipage}[t]{0.48\linewidth}\centering
\begin{tikzpicture}
  \tkzTabInit[lgt=1.5,espcl=2.2]{$x$/1,$\cos(x)$/2}{$-\pi$,$0$,$\pi$}
  \tkzTabVar{-/$-1$,+/$1$,-/$-1$}
\end{tikzpicture}
\end{minipage}\hfill
\begin{minipage}[t]{0.48\linewidth}\centering
\begin{tikzpicture}
  \tkzTabInit[lgt=1.5,espcl=1.5]{$x$/1.3,$\sin(x)$/2}{$-\pi$,$-\dfrac{\pi}{2}$,$\dfrac{\pi}{2}$,$\pi$}
  \tkzTabVar{+/$0$,-/$-1$,+/$1$,-/$0$}
\end{tikzpicture}
\end{minipage}
\end{proprietebox}

\begin{methodebox}[Méthode 12 --- Dresser le tableau de variation sur un autre intervalle]
\etapeexemples
Dresser le tableau de variation, à l'aide des courbes de la propriété 7 :
\begin{questions}[label=\alph*)]
  \item de la fonction cosinus sur $\intff{\pi}{3\pi}$ ;
  \item de la fonction sinus sur $\intff{\dfrac{\pi}{2}}{\dfrac{5\pi}{2}}$ ;
  \item de la fonction cosinus sur $\intff{-3\pi}{-\pi}$.
\end{questions}
\begin{correction}
\cas{a)}\par\nopagebreak
\begin{minipage}[c]{0.46\linewidth}\raggedright
La courbe part de $-1$, \res{monte à $1$ en $x=2\pi$, puis redescend à $-1$}.
\end{minipage}\hfill
\begin{minipage}[c]{0.50\linewidth}\centering
\begin{tikzpicture}
  \tkzTabInit[lgt=1.5,espcl=2]{$x$/1,$\cos(x)$/2}{$\pi$,$2\pi$,$3\pi$}
  \tkzTabVar{-/$-1$,+/$1$,-/$-1$}
\end{tikzpicture}
\end{minipage}
\cas{b)}\par\nopagebreak
\begin{minipage}[c]{0.46\linewidth}\raggedright
La courbe part de $1$, \res{descend à $-1$ en $x=\dfrac{3\pi}{2}$, puis remonte à $1$}.
\end{minipage}\hfill
\begin{minipage}[c]{0.50\linewidth}\centering
\begin{tikzpicture}
  \tkzTabInit[lgt=1.5,espcl=2]{$x$/1.3,$\sin(x)$/2}{$\dfrac{\pi}{2}$,$\dfrac{3\pi}{2}$,$\dfrac{5\pi}{2}$}
  \tkzTabVar{+/$1$,-/$-1$,+/$1$}
\end{tikzpicture}
\end{minipage}
\cas{c)}\par\nopagebreak
\begin{minipage}[c]{0.46\linewidth}\raggedright
La période est $2\pi$ : la courbe refait sur $\intff{-3\pi}{-\pi}$ ce qu'elle fait sur
$\intff{-\pi}{\pi}$. On retire $2\pi$ aux abscisses du tableau de la propriété 7 :
\res{croissante de $-1$ à $1$ sur $\intff{-3\pi}{-2\pi}$, puis décroissante}.
\end{minipage}\hfill
\begin{minipage}[c]{0.50\linewidth}\centering
\begin{tikzpicture}
  \tkzTabInit[lgt=1.5,espcl=2]{$x$/1,$\cos(x)$/2}{$-3\pi$,$-2\pi$,$-\pi$}
  \tkzTabVar{-/$-1$,+/$1$,-/$-1$}
\end{tikzpicture}
\end{minipage}
\end{correction}

\etape{À vous de jouer}
Dresser le tableau de variation :
\begin{questions}[label=\alph*)]
  \item de la fonction cosinus sur $\intff{0}{2\pi}$ ;
  \item de la fonction sinus sur $\intff{-\dfrac{3\pi}{2}}{\dfrac{\pi}{2}}$ ;
  \item de la fonction cosinus sur $\intff{2\pi}{4\pi}$.
\end{questions}
\reponse{a)~décroissante de $1$ à $-1$ sur $\intff{0}{\pi}$, croissante de $-1$ à $1$ sur $\intff{\pi}{2\pi}$ ;\;
b)~décroissante de $1$ à $-1$ sur $\intff{-\dfrac{3\pi}{2}}{-\dfrac{\pi}{2}}$, croissante de $-1$ à $1$ sur $\intff{-\dfrac{\pi}{2}}{\dfrac{\pi}{2}}$ ;\;
c)~décroissante de $1$ à $-1$ sur $\intff{2\pi}{3\pi}$, croissante de $-1$ à $1$ sur $\intff{3\pi}{4\pi}$.}
\end{methodebox}

\begin{methodebox}[Méthode 13 --- Démontrer qu'une fonction est périodique]
\textbf{But :} montrer que $f(x+T)=f(x)$ pour tout nombre réel $x$. Pour cela, on remplace
$x$ par $x+T$ dans la fonction, puis on fait apparaître un tour, $2\pi$.
\etapeexemples
Ces fonctions sont définies sur $\R$.
\begin{questions}[label=\alph*)]
  \item Démontrer que $f(x)=3\cos(x)+1$ est périodique de période $2\pi$.
  \item Démontrer que $g(x)=\sin(2x)$ est périodique de période $\pi$.
  \item Démontrer que $h(x)=\cos\left(\dfrac{x}{3}\right)$ est périodique de période $6\pi$.
\end{questions}
\begin{correction}
\cas{a)}
\begin{calculs}
f(x+2\pi) &= 3\cos(x+2\pi)+1 &\qquad& \rem{on remplace $x$ par $x+2\pi$}\\
f(x+2\pi) &= 3\cos(x)+1 && \rem{$\cos(x+2\pi)=\cos(x)$ : un tour de plus}\\
f(x+2\pi) &= f(x) && \rem{on reconnaît $f(x)$}
\end{calculs}
\explique{La fonction $f$ est \res{périodique de période $2\pi$}.}{on conclut par une phrase}
\cas{b)}
\begin{calculs}
g(x+\pi) &= \sin\big(2(x+\pi)\big) &\qquad& \rem{on remplace $x$ par $x+\pi$}\\
g(x+\pi) &= \sin(2x+2\pi) && \rem{on développe : un tour apparaît}\\
g(x+\pi) &= \sin(2x) && \rem{$\sin(X+2\pi)=\sin(X)$, avec $X=2x$}\\
g(x+\pi) &= g(x) && \rem{on reconnaît $g(x)$}
\end{calculs}
\explique{La fonction $g$ est \res{périodique de période $\pi$}.}{sa courbe se répète deux fois plus souvent que celle du sinus}
\cas{c)}
\begin{calculs}
h(x+6\pi) &= \cos\left(\dfrac{x+6\pi}{3}\right) &\qquad& \rem{on remplace $x$ par $x+6\pi$}\\
h(x+6\pi) &= \cos\left(\dfrac{x}{3}+2\pi\right) && \rem{$\tfrac{6\pi}{3}$ vaut $2\pi$}\\
h(x+6\pi) &= \cos\left(\dfrac{x}{3}\right) && \rem{un tour de plus}\\
h(x+6\pi) &= h(x) && \rem{on reconnaît $h(x)$}
\end{calculs}
\explique{La fonction $h$ est \res{périodique de période $6\pi$}.}{on conclut par une phrase}
\end{correction}

\etape{À vous de jouer}
Ces fonctions sont définies sur $\R$.
\begin{questions}[label=\alph*)]
  \item Démontrer que $k(x)=2-\sin(x)$ est périodique de période $2\pi$.
  \item Démontrer que $m(x)=\cos(4x)$ est périodique de période $\dfrac{\pi}{2}$.
  \item Démontrer que $n(x)=\sin\left(\dfrac{x}{2}\right)$ est périodique de période $4\pi$.
\end{questions}
\reponse{a)~$k(x+2\pi)=k(x)$ ;\; b)~$m\left(x+\dfrac{\pi}{2}\right)=m(x)$ ;\; c)~$n(x+4\pi)=n(x)$.}
\end{methodebox}

\afaire{exercices 35 à 38}

\partiecours{Résolution graphique d'équations et d'inéquations}

Comme pour toute fonction, on lit les solutions sur la courbe : les abscisses de ses points
d'intersection avec une droite horizontale. La courbe montre d'un seul coup toutes les
solutions d'un intervalle.

\begin{methodebox}[Méthode 14 --- Résoudre graphiquement une équation ou une inéquation]
\etapeexemples
Un carreau vaut $\dfrac{\pi}{6}$ en abscisse et $\dfrac{1}{2}$ en ordonnée. Résoudre graphiquement chaque équation, puis chaque inéquation :
\begin{questions}[label=\alph*)]
  \item $\sin(x)=\dfrac{1}{2}$ sur $\intff{-2\pi}{2\pi}$ ;
  \item $\cos(x)=-\dfrac{1}{2}$ sur $\intff{-2\pi}{2\pi}$ ;
  \item $\cos(x)=\dfrac{\sqrt{3}}{2}$ sur $\intff{0}{2\pi}$ ;
  \item $\sin(x)\ge\dfrac{1}{2}$ sur $\intff{0}{2\pi}$ ;
  \item $\cos(x)<\dfrac{1}{2}$ sur $\intff{-\pi}{\pi}$ ;
  \item $\sin(x)\le-\dfrac{1}{2}$ sur $\intff{-2\pi}{2\pi}$.
\end{questions}
\begin{center}
\begin{tikzpicture}[x=1.2cm,y=1.2cm]
  \reperetrigo{-12}{12}\courbecos{-12}{12}\gradcos
  \node[coulCourbe,font=\footnotesize,fill=white,inner sep=1pt] at (2.2,1.25) {$y=\cos(x)$};
\end{tikzpicture}

\medskip
\begin{tikzpicture}[x=1.2cm,y=1.2cm]
  \reperetrigo{-12}{12}\courbesin{-12}{12}\gradsin
  \node[coulCourbe,font=\footnotesize,fill=white,inner sep=1pt] at (3.3,1.25) {$y=\sin(x)$};
\end{tikzpicture}
\end{center}
\begin{correction}
\cas{a)}\par\nopagebreak
\explique{On trace la droite d'équation $y=\dfrac{1}{2}$ et on lit les abscisses des points d'intersection.}{la droite coupe la courbe du sinus en quatre points}
\explique{$S=\res{\ens{-\dfrac{11\pi}{6}\,;-\dfrac{7\pi}{6}\,;\dfrac{\pi}{6}\,;\dfrac{5\pi}{6}}}$.}{$\tfrac{\pi}{6}$ : un carreau à droite de $0$}
\begin{center}
\begin{tikzpicture}[x=1.2cm,y=1.2cm]
  \reperetrigo{-12}{12}\courbesin{-12}{12}
  \droiteh{-12}{12}{0.5}{$y=\frac{1}{2}$}
  \lirex{-11}{0.5}\lirex{-7}{0.5}\lirex{1}{0.5}\lirex{5}{0.5}
  \gradpi{-11/-\frac{11\pi}{6},-7/-\frac{7\pi}{6},1/\frac{\pi}{6},5/\frac{5\pi}{6}}
\end{tikzpicture}
\end{center}
\cas{b)}\par\nopagebreak
\explique{On trace la droite d'équation $y=-\dfrac{1}{2}$ et on lit les abscisses des points d'intersection.}{même démarche, sur la courbe du cosinus}
\explique{$S=\res{\ens{-\dfrac{4\pi}{3}\,;-\dfrac{2\pi}{3}\,;\dfrac{2\pi}{3}\,;\dfrac{4\pi}{3}}}$.}{$\tfrac{2\pi}{3}$ : quatre carreaux à droite de $0$}
\begin{center}
\begin{tikzpicture}[x=1.2cm,y=1.2cm]
  \reperetrigo{-12}{12}\courbecos{-12}{12}
  \droiteh{-12}{12}{-0.5}{$y=-\frac{1}{2}$}
  \lirex{-8}{-0.5}\lirex{-4}{-0.5}\lirex{4}{-0.5}\lirex{8}{-0.5}
  \gradpi[above]{-8/-\frac{4\pi}{3},-4/-\frac{2\pi}{3},4/\frac{2\pi}{3},8/\frac{4\pi}{3}}
\end{tikzpicture}
\end{center}
\cas{c)}\par\nopagebreak
\begin{minipage}[c]{0.42\linewidth}\raggedright
\explique{$S=\res{\ens{\dfrac{\pi}{6}\,;\dfrac{11\pi}{6}}}$.}{seul $\intff{0}{2\pi}$ compte}
\end{minipage}\hfill
\begin{minipage}[c]{0.56\linewidth}\centering
\begin{tikzpicture}[x=1.2cm,y=1.2cm]
  \reperetrigo{0}{12}\courbecos{0}{12}
  \droiteh{0}{12}{0.866}{$y=\frac{\sqrt{3}}{2}$}
  \lirex{1}{0.866}\lirex{11}{0.866}
  \gradpi{1/\frac{\pi}{6},11/\frac{11\pi}{6}}
\end{tikzpicture}
\end{minipage}
\cas{d)}\par\nopagebreak
\begin{minipage}[c]{0.42\linewidth}\raggedright
\explique{$S=\res{\intff{\dfrac{\pi}{6}}{\dfrac{5\pi}{6}}}$.}{la courbe est sur ou au-dessus de la droite ; crochets fermés à cause de $\ge$}
\end{minipage}\hfill
\begin{minipage}[c]{0.56\linewidth}\centering
\begin{tikzpicture}[x=1.2cm,y=1.2cm]
  \reperetrigo{0}{12}\courbesin{0}{12}
  \droiteh{0}{12}{0.5}{$y=\frac{1}{2}$}
  \interaxe{1}{5}
  \lirex{1}{0.5}\lirex{5}{0.5}
  \gradpi{1/\frac{\pi}{6},5/\frac{5\pi}{6},12/2\pi}
\end{tikzpicture}
\end{minipage}
\cas{e)}\par\nopagebreak
\begin{minipage}[c]{0.42\linewidth}\raggedright
\explique{$S=\res{\intfo{-\pi}{-\dfrac{\pi}{3}}\cup\intof{\dfrac{\pi}{3}}{\pi}}$.}{la courbe est au-dessous de la droite, en deux morceaux ; crochets ouverts en $\pm\tfrac{\pi}{3}$ à cause de $<$}
\end{minipage}\hfill
\begin{minipage}[c]{0.56\linewidth}\centering
\begin{tikzpicture}[x=1.2cm,y=1.2cm]
  \reperetrigo{-6}{6}\courbecos{-6}{6}
  \droiteh{-6}{6}{0.5}{$y=\frac{1}{2}$}
  \interaxe{-6}{-2}\interaxe[coulCourbeC]{2}{6}
  \lirex{-2}{0.5}\lirex{2}{0.5}
  \gradpi{-6/-\pi,-2/-\frac{\pi}{3},2/\frac{\pi}{3},6/\pi}
\end{tikzpicture}
\end{minipage}
\cas{f)}\par\nopagebreak
\explique{$S=\res{\intff{-\dfrac{5\pi}{6}}{-\dfrac{\pi}{6}}\cup\intff{\dfrac{7\pi}{6}}{\dfrac{11\pi}{6}}}$.}{la courbe est sur ou au-dessous de la droite}
\begin{center}
\begin{tikzpicture}[x=1.2cm,y=1.2cm]
  \reperetrigo{-12}{12}\courbesin{-12}{12}
  \droiteh{-12}{12}{-0.5}{$y=-\frac{1}{2}$}
  \interaxe{-5}{-1}\interaxe[coulCourbeC]{7}{11}
  \lirex{-5}{-0.5}\lirex{-1}{-0.5}\lirex{7}{-0.5}\lirex{11}{-0.5}
  \gradpi[above]{-5/-\frac{5\pi}{6},-1/-\frac{\pi}{6},7/\frac{7\pi}{6},11/\frac{11\pi}{6}}
\end{tikzpicture}
\end{center}
\end{correction}

\etape{À vous de jouer}
Résoudre graphiquement, à l'aide des courbes ci-dessus :
\begin{center}
a)~$\sin(x)=-\dfrac{1}{2}$ sur $\intff{-2\pi}{2\pi}$ \qquad
b)~$\cos(x)=\dfrac{1}{2}$ sur $\intff{-2\pi}{2\pi}$ \qquad
c)~$\sin(x)=\dfrac{\sqrt{3}}{2}$ sur $\intff{0}{2\pi}$\\[4pt]
d)~$\sin(x)\ge\dfrac{\sqrt{3}}{2}$ sur $\intff{0}{2\pi}$ \qquad
e)~$\cos(x)<-\dfrac{1}{2}$ sur $\intff{-\pi}{\pi}$ \qquad
f)~$\sin(x)\ge\dfrac{1}{2}$ sur $\intff{-2\pi}{2\pi}$
\end{center}
\reponse{a)~$\ens{-\dfrac{5\pi}{6}\,;-\dfrac{\pi}{6}\,;\dfrac{7\pi}{6}\,;\dfrac{11\pi}{6}}$ ;\;
b)~$\ens{-\dfrac{5\pi}{3}\,;-\dfrac{\pi}{3}\,;\dfrac{\pi}{3}\,;\dfrac{5\pi}{3}}$ ;\;
c)~$\ens{\dfrac{\pi}{3}\,;\dfrac{2\pi}{3}}$ ;\;
d)~$\intff{\dfrac{\pi}{3}}{\dfrac{2\pi}{3}}$ ;\;
e)~$\intfo{-\pi}{-\dfrac{2\pi}{3}}\cup\intof{\dfrac{2\pi}{3}}{\pi}$ ;\;
f)~$\intff{-\dfrac{11\pi}{6}}{-\dfrac{7\pi}{6}}\cup\intff{\dfrac{\pi}{6}}{\dfrac{5\pi}{6}}$.}
\end{methodebox}

\afaire{exercices 39 à 42}

\end{document}
